% concatenated from main.tex, math_commands.tex
\documentclass{article}
%%%%% NEW MATH DEFINITIONS %%%%%

\usepackage{amsmath,amsfonts,bm}

% Mark sections of captions for referring to divisions of figures
\newcommand{\figleft}{{\em (Left)}}
\newcommand{\figcenter}{{\em (Center)}}
\newcommand{\figright}{{\em (Right)}}
\newcommand{\figtop}{{\em (Top)}}
\newcommand{\figbottom}{{\em (Bottom)}}
\newcommand{\captiona}{{\em (a)}}
\newcommand{\captionb}{{\em (b)}}
\newcommand{\captionc}{{\em (c)}}
\newcommand{\captiond}{{\em (d)}}

% Highlight a newly defined term
\newcommand{\newterm}[1]{{\bf #1}}


% Figure reference, lower-case.
\def\figref#1{figure~\ref{#1}}
% Figure reference, capital. For start of sentence
\def\Figref#1{Figure~\ref{#1}}
\def\twofigref#1#2{figures \ref{#1} and \ref{#2}}
\def\quadfigref#1#2#3#4{figures \ref{#1}, \ref{#2}, \ref{#3} and \ref{#4}}
% Section reference, lower-case.
\def\secref#1{section~\ref{#1}}
% Section reference, capital.
\def\Secref#1{Section~\ref{#1}}
% Reference to two sections.
\def\twosecrefs#1#2{sections \ref{#1} and \ref{#2}}
% Reference to three sections.
\def\secrefs#1#2#3{sections \ref{#1}, \ref{#2} and \ref{#3}}
% Reference to an equation, lower-case.
\def\eqref#1{equation~\ref{#1}}
% Reference to an equation, upper case
\def\Eqref#1{Equation~\ref{#1}}
% A raw reference to an equation---avoid using if possible
\def\plaineqref#1{\ref{#1}}
% Reference to a chapter, lower-case.
\def\chapref#1{chapter~\ref{#1}}
% Reference to an equation, upper case.
\def\Chapref#1{Chapter~\ref{#1}}
% Reference to a range of chapters
\def\rangechapref#1#2{chapters\ref{#1}--\ref{#2}}
% Reference to an algorithm, lower-case.
\def\algref#1{algorithm~\ref{#1}}
% Reference to an algorithm, upper case.
\def\Algref#1{Algorithm~\ref{#1}}
\def\twoalgref#1#2{algorithms \ref{#1} and \ref{#2}}
\def\Twoalgref#1#2{Algorithms \ref{#1} and \ref{#2}}
% Reference to a part, lower case
\def\partref#1{part~\ref{#1}}
% Reference to a part, upper case
\def\Partref#1{Part~\ref{#1}}
\def\twopartref#1#2{parts \ref{#1} and \ref{#2}}

\def\ceil#1{\lceil #1 \rceil}
\def\floor#1{\lfloor #1 \rfloor}
\def\1{\bm{1}}
\newcommand{\train}{\mathcal{D}}
\newcommand{\valid}{\mathcal{D_{\mathrm{valid}}}}
\newcommand{\test}{\mathcal{D_{\mathrm{test}}}}

\def\eps{{\epsilon}}


% Random variables
\def\reta{{\textnormal{$\eta$}}}
\def\ra{{\textnormal{a}}}
\def\rb{{\textnormal{b}}}
\def\rc{{\textnormal{c}}}
\def\rd{{\textnormal{d}}}
\def\re{{\textnormal{e}}}
\def\rf{{\textnormal{f}}}
\def\rg{{\textnormal{g}}}
\def\rh{{\textnormal{h}}}
\def\ri{{\textnormal{i}}}
\def\rj{{\textnormal{j}}}
\def\rk{{\textnormal{k}}}
\def\rl{{\textnormal{l}}}
% rm is already a command, just don't name any random variables m
\def\rn{{\textnormal{n}}}
\def\ro{{\textnormal{o}}}
\def\rp{{\textnormal{p}}}
\def\rq{{\textnormal{q}}}
\def\rr{{\textnormal{r}}}
\def\rs{{\textnormal{s}}}
\def\rt{{\textnormal{t}}}
\def\ru{{\textnormal{u}}}
\def\rv{{\textnormal{v}}}
\def\rw{{\textnormal{w}}}
\def\rx{{\textnormal{x}}}
\def\ry{{\textnormal{y}}}
\def\rz{{\textnormal{z}}}

% Random vectors
\def\rvepsilon{{\mathbf{\epsilon}}}
\def\rvtheta{{\mathbf{\theta}}}
\def\rva{{\mathbf{a}}}
\def\rvb{{\mathbf{b}}}
\def\rvc{{\mathbf{c}}}
\def\rvd{{\mathbf{d}}}
\def\rve{{\mathbf{e}}}
\def\rvf{{\mathbf{f}}}
\def\rvg{{\mathbf{g}}}
\def\rvh{{\mathbf{h}}}
\def\rvu{{\mathbf{i}}}
\def\rvj{{\mathbf{j}}}
\def\rvk{{\mathbf{k}}}
\def\rvl{{\mathbf{l}}}
\def\rvm{{\mathbf{m}}}
\def\rvn{{\mathbf{n}}}
\def\rvo{{\mathbf{o}}}
\def\rvp{{\mathbf{p}}}
\def\rvq{{\mathbf{q}}}
\def\rvr{{\mathbf{r}}}
\def\rvs{{\mathbf{s}}}
\def\rvt{{\mathbf{t}}}
\def\rvu{{\mathbf{u}}}
\def\rvv{{\mathbf{v}}}
\def\rvw{{\mathbf{w}}}
\def\rvx{{\mathbf{x}}}
\def\rvy{{\mathbf{y}}}
\def\rvz{{\mathbf{z}}}

% Elements of random vectors
\def\erva{{\textnormal{a}}}
\def\ervb{{\textnormal{b}}}
\def\ervc{{\textnormal{c}}}
\def\ervd{{\textnormal{d}}}
\def\erve{{\textnormal{e}}}
\def\ervf{{\textnormal{f}}}
\def\ervg{{\textnormal{g}}}
\def\ervh{{\textnormal{h}}}
\def\ervi{{\textnormal{i}}}
\def\ervj{{\textnormal{j}}}
\def\ervk{{\textnormal{k}}}
\def\ervl{{\textnormal{l}}}
\def\ervm{{\textnormal{m}}}
\def\ervn{{\textnormal{n}}}
\def\ervo{{\textnormal{o}}}
\def\ervp{{\textnormal{p}}}
\def\ervq{{\textnormal{q}}}
\def\ervr{{\textnormal{r}}}
\def\ervs{{\textnormal{s}}}
\def\ervt{{\textnormal{t}}}
\def\ervu{{\textnormal{u}}}
\def\ervv{{\textnormal{v}}}
\def\ervw{{\textnormal{w}}}
\def\ervx{{\textnormal{x}}}
\def\ervy{{\textnormal{y}}}
\def\ervz{{\textnormal{z}}}

% Random matrices
\def\rmA{{\mathbf{A}}}
\def\rmB{{\mathbf{B}}}
\def\rmC{{\mathbf{C}}}
\def\rmD{{\mathbf{D}}}
\def\rmE{{\mathbf{E}}}
\def\rmF{{\mathbf{F}}}
\def\rmG{{\mathbf{G}}}
\def\rmH{{\mathbf{H}}}
\def\rmI{{\mathbf{I}}}
\def\rmJ{{\mathbf{J}}}
\def\rmK{{\mathbf{K}}}
\def\rmL{{\mathbf{L}}}
\def\rmM{{\mathbf{M}}}
\def\rmN{{\mathbf{N}}}
\def\rmO{{\mathbf{O}}}
\def\rmP{{\mathbf{P}}}
\def\rmQ{{\mathbf{Q}}}
\def\rmR{{\mathbf{R}}}
\def\rmS{{\mathbf{S}}}
\def\rmT{{\mathbf{T}}}
\def\rmU{{\mathbf{U}}}
\def\rmV{{\mathbf{V}}}
\def\rmW{{\mathbf{W}}}
\def\rmX{{\mathbf{X}}}
\def\rmY{{\mathbf{Y}}}
\def\rmZ{{\mathbf{Z}}}

% Elements of random matrices
\def\ermA{{\textnormal{A}}}
\def\ermB{{\textnormal{B}}}
\def\ermC{{\textnormal{C}}}
\def\ermD{{\textnormal{D}}}
\def\ermE{{\textnormal{E}}}
\def\ermF{{\textnormal{F}}}
\def\ermG{{\textnormal{G}}}
\def\ermH{{\textnormal{H}}}
\def\ermI{{\textnormal{I}}}
\def\ermJ{{\textnormal{J}}}
\def\ermK{{\textnormal{K}}}
\def\ermL{{\textnormal{L}}}
\def\ermM{{\textnormal{M}}}
\def\ermN{{\textnormal{N}}}
\def\ermO{{\textnormal{O}}}
\def\ermP{{\textnormal{P}}}
\def\ermQ{{\textnormal{Q}}}
\def\ermR{{\textnormal{R}}}
\def\ermS{{\textnormal{S}}}
\def\ermT{{\textnormal{T}}}
\def\ermU{{\textnormal{U}}}
\def\ermV{{\textnormal{V}}}
\def\ermW{{\textnormal{W}}}
\def\ermX{{\textnormal{X}}}
\def\ermY{{\textnormal{Y}}}
\def\ermZ{{\textnormal{Z}}}

% Vectors
\def\vzero{{\bm{0}}}
\def\vone{{\bm{1}}}
\def\vmu{{\bm{\mu}}}
\def\vtheta{{\bm{\theta}}}
\def\va{{\bm{a}}}
\def\vb{{\bm{b}}}
\def\vc{{\bm{c}}}
\def\vd{{\bm{d}}}
\def\ve{{\bm{e}}}
\def\vf{{\bm{f}}}
\def\vg{{\bm{g}}}
\def\vh{{\bm{h}}}
\def\vi{{\bm{i}}}
\def\vj{{\bm{j}}}
\def\vk{{\bm{k}}}
\def\vl{{\bm{l}}}
\def\vm{{\bm{m}}}
\def\vn{{\bm{n}}}
\def\vo{{\bm{o}}}
\def\vp{{\bm{p}}}
\def\vq{{\bm{q}}}
\def\vr{{\bm{r}}}
\def\vs{{\bm{s}}}
\def\vt{{\bm{t}}}
\def\vu{{\bm{u}}}
\def\vv{{\bm{v}}}
\def\vw{{\bm{w}}}
\def\vx{{\bm{x}}}
\def\vy{{\bm{y}}}
\def\vz{{\bm{z}}}

% Elements of vectors
\def\evalpha{{\alpha}}
\def\evbeta{{\beta}}
\def\evepsilon{{\epsilon}}
\def\evlambda{{\lambda}}
\def\evomega{{\omega}}
\def\evmu{{\mu}}
\def\evpsi{{\psi}}
\def\evsigma{{\sigma}}
\def\evtheta{{\theta}}
\def\eva{{a}}
\def\evb{{b}}
\def\evc{{c}}
\def\evd{{d}}
\def\eve{{e}}
\def\evf{{f}}
\def\evg{{g}}
\def\evh{{h}}
\def\evi{{i}}
\def\evj{{j}}
\def\evk{{k}}
\def\evl{{l}}
\def\evm{{m}}
\def\evn{{n}}
\def\evo{{o}}
\def\evp{{p}}
\def\evq{{q}}
\def\evr{{r}}
\def\evs{{s}}
\def\evt{{t}}
\def\evu{{u}}
\def\evv{{v}}
\def\evw{{w}}
\def\evx{{x}}
\def\evy{{y}}
\def\evz{{z}}

% Matrix
\def\mA{{\bm{A}}}
\def\mB{{\bm{B}}}
\def\mC{{\bm{C}}}
\def\mD{{\bm{D}}}
\def\mE{{\bm{E}}}
\def\mF{{\bm{F}}}
\def\mG{{\bm{G}}}
\def\mH{{\bm{H}}}
\def\mI{{\bm{I}}}
\def\mJ{{\bm{J}}}
\def\mK{{\bm{K}}}
\def\mL{{\bm{L}}}
\def\mM{{\bm{M}}}
\def\mN{{\bm{N}}}
\def\mO{{\bm{O}}}
\def\mP{{\bm{P}}}
\def\mQ{{\bm{Q}}}
\def\mR{{\bm{R}}}
\def\mS{{\bm{S}}}
\def\mT{{\bm{T}}}
\def\mU{{\bm{U}}}
\def\mV{{\bm{V}}}
\def\mW{{\bm{W}}}
\def\mX{{\bm{X}}}
\def\mY{{\bm{Y}}}
\def\mZ{{\bm{Z}}}
\def\mBeta{{\bm{\beta}}}
\def\mPhi{{\bm{\Phi}}}
\def\mLambda{{\bm{\Lambda}}}
\def\mSigma{{\bm{\Sigma}}}

% Tensor
\DeclareMathAlphabet{\mathsfit}{\encodingdefault}{\sfdefault}{m}{sl}
\SetMathAlphabet{\mathsfit}{bold}{\encodingdefault}{\sfdefault}{bx}{n}
\newcommand{\tens}[1]{\bm{\mathsfit{#1}}}
\def\tA{{\tens{A}}}
\def\tB{{\tens{B}}}
\def\tC{{\tens{C}}}
\def\tD{{\tens{D}}}
\def\tE{{\tens{E}}}
\def\tF{{\tens{F}}}
\def\tG{{\tens{G}}}
\def\tH{{\tens{H}}}
\def\tI{{\tens{I}}}
\def\tJ{{\tens{J}}}
\def\tK{{\tens{K}}}
\def\tL{{\tens{L}}}
\def\tM{{\tens{M}}}
\def\tN{{\tens{N}}}
\def\tO{{\tens{O}}}
\def\tP{{\tens{P}}}
\def\tQ{{\tens{Q}}}
\def\tR{{\tens{R}}}
\def\tS{{\tens{S}}}
\def\tT{{\tens{T}}}
\def\tU{{\tens{U}}}
\def\tV{{\tens{V}}}
\def\tW{{\tens{W}}}
\def\tX{{\tens{X}}}
\def\tY{{\tens{Y}}}
\def\tZ{{\tens{Z}}}


% Graph
\def\gA{{\mathcal{A}}}
\def\gB{{\mathcal{B}}}
\def\gC{{\mathcal{C}}}
\def\gD{{\mathcal{D}}}
\def\gE{{\mathcal{E}}}
\def\gF{{\mathcal{F}}}
\def\gG{{\mathcal{G}}}
\def\gH{{\mathcal{H}}}
\def\gI{{\mathcal{I}}}
\def\gJ{{\mathcal{J}}}
\def\gK{{\mathcal{K}}}
\def\gL{{\mathcal{L}}}
\def\gM{{\mathcal{M}}}
\def\gN{{\mathcal{N}}}
\def\gO{{\mathcal{O}}}
\def\gP{{\mathcal{P}}}
\def\gQ{{\mathcal{Q}}}
\def\gR{{\mathcal{R}}}
\def\gS{{\mathcal{S}}}
\def\gT{{\mathcal{T}}}
\def\gU{{\mathcal{U}}}
\def\gV{{\mathcal{V}}}
\def\gW{{\mathcal{W}}}
\def\gX{{\mathcal{X}}}
\def\gY{{\mathcal{Y}}}
\def\gZ{{\mathcal{Z}}}

% Sets
\def\sA{{\mathbb{A}}}
\def\sB{{\mathbb{B}}}
\def\sC{{\mathbb{C}}}
\def\sD{{\mathbb{D}}}
% Don't use a set called E, because this would be the same as our symbol
% for expectation.
\def\sF{{\mathbb{F}}}
\def\sG{{\mathbb{G}}}
\def\sH{{\mathbb{H}}}
\def\sI{{\mathbb{I}}}
\def\sJ{{\mathbb{J}}}
\def\sK{{\mathbb{K}}}
\def\sL{{\mathbb{L}}}
\def\sM{{\mathbb{M}}}
\def\sN{{\mathbb{N}}}
\def\sO{{\mathbb{O}}}
\def\sP{{\mathbb{P}}}
\def\sQ{{\mathbb{Q}}}
\def\sR{{\mathbb{R}}}
\def\sS{{\mathbb{S}}}
\def\sT{{\mathbb{T}}}
\def\sU{{\mathbb{U}}}
\def\sV{{\mathbb{V}}}
\def\sW{{\mathbb{W}}}
\def\sX{{\mathbb{X}}}
\def\sY{{\mathbb{Y}}}
\def\sZ{{\mathbb{Z}}}

% Entries of a matrix
\def\emLambda{{\Lambda}}
\def\emA{{A}}
\def\emB{{B}}
\def\emC{{C}}
\def\emD{{D}}
\def\emE{{E}}
\def\emF{{F}}
\def\emG{{G}}
\def\emH{{H}}
\def\emI{{I}}
\def\emJ{{J}}
\def\emK{{K}}
\def\emL{{L}}
\def\emM{{M}}
\def\emN{{N}}
\def\emO{{O}}
\def\emP{{P}}
\def\emQ{{Q}}
\def\emR{{R}}
\def\emS{{S}}
\def\emT{{T}}
\def\emU{{U}}
\def\emV{{V}}
\def\emW{{W}}
\def\emX{{X}}
\def\emY{{Y}}
\def\emZ{{Z}}
\def\emSigma{{\Sigma}}

% entries of a tensor
% Same font as tensor, without \bm wrapper
\newcommand{\etens}[1]{\mathsfit{#1}}
\def\etLambda{{\etens{\Lambda}}}
\def\etA{{\etens{A}}}
\def\etB{{\etens{B}}}
\def\etC{{\etens{C}}}
\def\etD{{\etens{D}}}
\def\etE{{\etens{E}}}
\def\etF{{\etens{F}}}
\def\etG{{\etens{G}}}
\def\etH{{\etens{H}}}
\def\etI{{\etens{I}}}
\def\etJ{{\etens{J}}}
\def\etK{{\etens{K}}}
\def\etL{{\etens{L}}}
\def\etM{{\etens{M}}}
\def\etN{{\etens{N}}}
\def\etO{{\etens{O}}}
\def\etP{{\etens{P}}}
\def\etQ{{\etens{Q}}}
\def\etR{{\etens{R}}}
\def\etS{{\etens{S}}}
\def\etT{{\etens{T}}}
\def\etU{{\etens{U}}}
\def\etV{{\etens{V}}}
\def\etW{{\etens{W}}}
\def\etX{{\etens{X}}}
\def\etY{{\etens{Y}}}
\def\etZ{{\etens{Z}}}

% The true underlying data generating distribution
\newcommand{\pdata}{p_{\rm{data}}}
% The empirical distribution defined by the training set
\newcommand{\ptrain}{\hat{p}_{\rm{data}}}
\newcommand{\Ptrain}{\hat{P}_{\rm{data}}}
% The model distribution
\newcommand{\pmodel}{p_{\rm{model}}}
\newcommand{\Pmodel}{P_{\rm{model}}}
\newcommand{\ptildemodel}{\tilde{p}_{\rm{model}}}
% Stochastic autoencoder distributions
\newcommand{\pencode}{p_{\rm{encoder}}}
\newcommand{\pdecode}{p_{\rm{decoder}}}
\newcommand{\precons}{p_{\rm{reconstruct}}}

\newcommand{\laplace}{\mathrm{Laplace}} % Laplace distribution

\newcommand{\E}{\mathbb{E}}
\newcommand{\Ls}{\mathcal{L}}
\newcommand{\R}{\mathbb{R}}
\newcommand{\emp}{\tilde{p}}
\newcommand{\lr}{\alpha}
\newcommand{\reg}{\lambda}
\newcommand{\rect}{\mathrm{rectifier}}
\newcommand{\softmax}{\mathrm{softmax}}
\newcommand{\sigmoid}{\sigma}
\newcommand{\softplus}{\zeta}
\newcommand{\KL}{D_{\mathrm{KL}}}
\newcommand{\Var}{\mathrm{Var}}
\newcommand{\standarderror}{\mathrm{SE}}
\newcommand{\Cov}{\mathrm{Cov}}
% Wolfram Mathworld says $L^2$ is for function spaces and $\ell^2$ is for vectors
% But then they seem to use $L^2$ for vectors throughout the site, and so does
% wikipedia.
\newcommand{\normlzero}{L^0}
\newcommand{\normlone}{L^1}
\newcommand{\normltwo}{L^2}
\newcommand{\normlp}{L^p}
\newcommand{\normmax}{L^\infty}

\newcommand{\parents}{Pa} % See usage in notation.tex. Chosen to match Daphne's book.

\DeclareMathOperator*{\argmax}{arg\,max}
\DeclareMathOperator*{\argmin}{arg\,min}

\DeclareMathOperator{\sign}{sign}
\DeclareMathOperator{\Tr}{Tr}
\let\ab\allowbreak


%%%%%%%%%%%%%% self-defined command %%%%%%%%%%
\usepackage{amsmath}
\usepackage{mathtools}
\usepackage{amsthm}
\usepackage{xcolor}

\newtheorem{theorem}{Theorem}[section]
\newtheorem{lemma}[theorem]{Lemma}
% \newtheorem{remark}[theorem]{Remark}
\newtheorem{assumption}[theorem]{Assumption}
% \newtheorem{definition}[theorem]{Definition}
\newtheorem{proposition}[theorem]{Proposition}
\newtheorem{corollary}[theorem]{Corollary}
\newtheorem{example}{Example}[section]

\newtheorem{definition}{Definition}[section]
\newtheorem{remark}{Remark}[section]


\newcommand{\Rbb}{\mathbb{R}}
\newcommand{\Ebb}[1]{\mathbb{E}\left[#1\right]}
\newcommand{\Vcal}{\mathcal{V}}
\newcommand{\Scal}{\mathcal{S}}
\newcommand{\Ical}{\mathcal{I}}
\newcommand{\Lcal}{\mathcal{L}}
\newcommand{\Zcal}{\mathcal{Z}}
\newcommand{\Xcal}{\mathcal{X}}
\newcommand{\Ncal}{\mathcal{N}}
\newcommand{\Fcal}{\mathcal{F}}


\newcommand{\zsf}{\mathsf{z}}
\newcommand{\xsf}{\mathsf{x}}

\newcommand{\vep}{\varepsilon}
\newcommand{\Pdim}{\text{Pdim}}
\newcommand{\Ocal}[1]{\mathcal{O}\left(#1\right)}

\usepackage{hyperref}
\usepackage{url}
\usepackage{amsmath,amssymb,mathtools} 
\usepackage{algorithm}                  
\usepackage{algpseudocode}              
\usepackage{float}                     
\usepackage{cleveref}   
\usepackage{multirow} 
\usepackage{PRIMEarxiv}
\usepackage{natbib}
\usepackage{amssymb}
\usepackage{amsmath}
\usepackage{array}
\usepackage[utf8]{inputenc} % allow utf-8 input
\usepackage[T1]{fontenc}    % use 8-bit T1 fonts
\usepackage{hyperref}       % hyperlinks
\usepackage{url}            % simple URL typesetting
\usepackage{booktabs}       % professional-quality tables
\usepackage{amsfonts}       % blackboard math symbols
\usepackage{nicefrac}       % compact symbols for 1/2, etc.
\usepackage{microtype}      % microtypography
\usepackage{lipsum}
\usepackage{fancyhdr}       % header
\usepackage{graphicx}       % graphics
\graphicspath{{media/}}     % organize your images and other figures under media/ folder
\usepackage{subcaption} 
\usepackage{epstopdf}
%Header
\pagestyle{fancy}
\thispagestyle{empty}
\rhead{ \textit{ }} 

\fancyhead[LO]{}

\newcommand{\fix}{\marginpar{FIX}}
\newcommand{\new}{\marginpar{NEW}}
\newcommand{\srankV}{r_{\text{s}}(V)}
% Define a norm command for convenience

\newcommand{\conv}{\operatorname{conv}}
\newcommand{\Per}{\mathrm{Per}}
\newcommand{\diam}{\mathrm{diam}}
\newcommand{\dist}{\operatorname{dist}}
% Specific norm notations
\newcommand{\opnormspace}{\mathcal{L}(H^{1/2}(\partial\Omega), H^{-1/2}(\partial\Omega))}
\newcommand{\opnormspaceprojected}{\mathcal{L}(V_N, W_N)}
\newcommand{\opnormglobal}[1]{\norm{#1}_{\opnormspace}} % Operator norm on the full spaces H^1/2 -> H^-1/2
\newcommand{\opnormprojected}[1]{\norm{#1}_{\mathcal{L}(V_N, W_N)}} % Operator norm on the finite-dimensional spaces V_N -> W_N
\newcommand{\hsnorm}[1]{\norm{#1}_{HS}} % Hilbert-Schmidt norm
\newcommand{\lnorm}[1]{\norm{#1}_{L^1(\Omega)}}
\newcommand{\hhalfnorm}[1]{\norm{#1}_{H^{1/2}(\partial\Omega)}}
\newcommand{\hminushalfnorm}[1]{\norm{#1}_{H^{-1/2}(\partial\Omega)}}


\def\DtNobs{\\Lambda_{\gamma}^{obs}}
\def\DtNgt{\Lambda_{\gamma}^{obs}}

\newcommand{\kc}[1]{\textcolor{red}{Ke: #1}}
\newcommand{\cy}[1]{\textcolor{blue}{Chugang: #1}}
\newcommand*\diff{\mathop{}\mathrm{d}}
  
%% Title
\title{Data Completion for Electrical Impedance Tomography by Conditional Diffusion Models}

\author{Ke Chen \\
Department of Mathematical Sciences\\
University of Delaware\\
\texttt{kechen@udel.edu}\\
\And
Haizhao Yang \\
Department of Mathematics\\
Department of Computer Science\\
University of Maryland, College Park \\
\texttt{hzyang@umd.edu} \\
\And
Chugang Yi\thanks{ Corresponding author. Authors are listed in alphabetical order.} \\
Department of Mathematics \\
University of Maryland, College Park \\
\texttt{chugang@umd.edu} \\
}

\begin{document}
\maketitle
\begin{abstract}
Data scarcity is a fundamental barrier in Electrical Impedance Tomography (EIT), as undersampled Dirichlet-to-Neumann (DtN) measurements can substantially degrade conductivity reconstructions. We address this bottleneck by completing partially observed DtN measurements using a diffusion-based generative model. Specifically, we train a conditional diffusion model to learn the distribution of DtN data and to infer full measurement vectors given partial observations. Our approach supports flexible source–receiver configurations and can be used as a plug-in preprocessing step with off-the-shelf EIT solvers. Under mild assumptions on the polygon conductivity class, we derive nonasymptotic end-to-end bounds on the distributional discrepancy between the completed and ground-truth DtN measurements. In numerical experiments, we couple the proposed diffusion  completion procedure with a deep learning–based inverse solver and compare its performance against the same solver with full measurement data. The results show that diffusion completion enables reconstructions that closely match the full-data baseline while using only $1\%$ of the measurements. In contrast, standard baselines such as matrix completion require $30\%$ of the measurements to achieve similar reconstruction quality.
\end{abstract}



\section{Introduction}

Inverse problems governed by partial differential equations (PDEs) arise across a broad range of scientific and engineering domains, including medical imaging \citep{arridge1999optical,cheney1990noser,borcea2002electrical}, geophysical exploration \citep{tarantola1982generalized,pratt1999seismic}, and materials characterization \citep{bonnet2005inverse,garcia2011non}. 
Many such inverse problems are nonlinear and severely ill-posed \citep{bal2012introduction, uhlmann2009electrical, isakov2006inverse}, and 
Electrical Impedance Tomography (EIT) is a canonical example of these difficulties \citep{alessandrini1988stable, allers1991stability}. In EIT, the goal is to determine the conductivity distribution within a medium by applying voltage patterns on its boundary and measuring the induced current responses. One of the key components of EIT is the Dirichlet-to-Neumann (DtN) operator, which maps voltage patterns to current responses. Theoretically, the conductivity can be uniquely determined from full knowledge of the DtN operator under mild conditions \citep{nachman1996global, astala2006calderon,uhlmann2013inverse, sylvester1987global}. 
In practice, the continuous DtN operator is often discretized into a \emph{DtN matrix} containing voltage-current measurements on a finite number of boundary nodes. 
Typical methods for solving EIT include iterative solvers with regularization \citep{hanke1997regularizing, haber2000optimization,jin2012reconstruction, vauhkonen2002tikhonov}, direct methods \citep{chow2021direct,chow2014direct, mueller2020d}, and recent deep learning approaches \citep{hamilton2018deep, ong2022integral, guo2023transformer, cen2023electrical, chen2024pseudo, molinaro2023neural, bhat2025neural}. However, when only limited measurements are available, as is common in medical imaging \citep{heines2023clinical, putensen2019electrical, scaramuzzo2024electrical}, reconstruction pipelines that presuppose fully sampled DtN measurements often fail to produce reliable solutions. The scarcity or incompleteness of measurements thus poses a central challenge in practical EIT.

To address a limited measurement budget, active acquisition strategies such as optimal experimental design (OED) and reinforcement learning (RL) aim to identify measurement configurations or sensing policies that extract the most informative measurements from a limited number of experiments. Their effectiveness has been demonstrated in EIT \citep{hyvonen2024bayesian,jin2024continuous} and other inverse problems \citep{huan2013simulation,shen2022learning,alexanderian2021optimal,jiang2024reinforced,go2025accurate,koval2025non}. 
Although these approaches can improve the quality of the collected data, they typically require many forward and inverse solves. Moreover, partial measurements can hinder the use of off-the-shelf reconstruction methods that expect complete measurements. For instance, in inverse scattering, limited aperture data leads to artifacts known as ``shadow regions''~\citep{colton2006target, li2015recovering}, while in EIT partial boundary data exacerbates the loss of spatial resolution and causes geometrical distortion~\citep{hauptmann2017approximation, alsaker2017direct}.


Another line of work seeks to reconstruct the full measurements from partial observations in many inverse problems \citep{harrach2015interpolation, hauptmann2017approximation,caubet2019new,dou2022data,bui2022bridging}, thereby bridging the gap to inverse solvers that presuppose complete measurements. For EIT, \citet{harrach2015interpolation} predict the missing voltages on current-driven boundary nodes in the difference Neumann-to-Dirichlet (NtD) matrix, relying on the geometry-dependent smoothness of difference NtD data; \citet{hauptmann2017approximation} approximate continuum NtD data from boundary measurements via a two-stage optimization that first lifts voltages to smooth boundary traces using Tikhonov regularization and then applies a prior-regularized least-squares fit to recover an approximate full NtD map; and the most closely related work, \citet{bui2022bridging}, completes unobserved entries in the off-diagonal blocks of the hierarchically partitioned DtN matrix by exploiting the low-rank structure inherited from the underlying elliptic operator \citep{bebendorf2003existence}. Specifically, \citet{bui2022bridging} achieve high-accuracy completion when the off-diagonal blocks of the DtN matrix are randomly sampled at a moderate sampling rate. 

These EIT data completion methods \citep{harrach2015interpolation,hauptmann2017approximation,bui2022bridging} leverage PDE-induced smoothness or low-rank structure, but may become less reliable under highly structured missingness, extremely sparse electrode layouts, or strong measurement noise. They typically perform completion in an instance-wise manner, requiring a problem-specific interpolation, regularization or matrix-completion routine to be executed for each new measurement set and acquisition configuration. In contrast, we propose a data-driven completion framework based on conditional diffusion models. By exploiting the low-dimensional structure of the conductivity distribution, our method amortizes completion across instances and enables accurate DtN completion under flexible and challenging measurement configurations. Diffusion models, also known as score-based generative models, have achieved remarkable success in learning complex distributions from high-dimensional data \citep{song2020score, ho2020denoising}, and recent theory shows that they effectively capture low-dimensional structure in high-dimensional distributions \citep{tang2024conditional,azangulov2024convergence, liang2025low}, a property that has been leveraged in recent EIT works \citep{alberti2022inverse, alberti2024manifold}. Building on the low-dimensional manifold structure of the conductivity distribution and the continuity of the forward map with respect to the conductivity, we show that the induced distribution of DtN measurements can be learned efficiently using a conditional diffusion model. The main contributions of this work can be summarized as follows: 
\begin{itemize}
  \item \textbf{DtN completion via conditional diffusion.} To the best of our knowledge, this is the first framework that learns the distribution of (discrete) DtN measurements via a diffusion model and performs DtN completion by posterior sampling conditioned on partial observations.

\item \textbf{Extreme undersampling performance.} Due to the ill-posedness of EIT, DtN measurements from different conductivities can be nearly indistinguishable, which hinders learning under heavy missingness. Our EIT-tailored normalization enhances contrast across measurements and stabilizes optimization of the neural network, enabling accurate completion from only $1\%$ of the measurements, a regime in which low-rank matrix completion fails.
\item \textbf{Flexible measurement configurations.} Whereas matrix completion relies on random missingness, our method accommodates flexible source–receiver patterns by encoding partial measurements as conditioning inputs to the score network. Thus our model can learn to infer missing entries under varied masks, providing robust completion under diverse measurement configurations.

  \item \textbf{Convergence guarantees.} By bridging diffusion models and EIT theory, we establish non-asymptotic end-to-end bounds for our DtN completion algorithm for polygonal conductivities.
  \item \textbf{Plug-and-play with inverse solvers.} The completed measurements can be used as a preprocessing step for off-the-shelf EIT solvers; when preceded by our completion step, a deep inverse network yields substantially better reconstructions than the same architecture trained directly on partial data.
\end{itemize}

The remainder of the paper is organized as follows. Section~\ref{sec:problem_setup} formalizes the problem; Section~\ref{sec:method} presents our diffusion completion framework for DtN measurements; Section~\ref{sec:theory} develops the error analysis; and Section~\ref{sec:numeric} reports numerical results.


\section{Problem Setup}\label{sec:problem_setup} 
Let $\Omega \subset \mathbb{R}^2$ be the unit disk, i.e., $\Omega = \{(x, y) \mid x^2 + y^2 < 1\} \subset \mathbb{R}^2$. The unknown parameter of interest is the conductivity $\gamma: \Omega \to \mathbb{R}^+$, where $\gamma(x)$ represents the conductivity at point $x \in \Omega$ and $\mathbb{R}^+$ denotes the set of positive real numbers. The forward problem is to solve the elliptic equation
\begin{equation}\label{eq:forward}
\nabla \cdot (\gamma(x)\nabla u(x)) = 0 \quad \text{in } \Omega, \quad u(x) = f(x) \quad \text{on the boundary} \, \partial \Omega\,,
\end{equation}
where $f: \partial \Omega \to \mathbb{R}$ is a prescribed voltage function on the boundary $\partial \Omega$. The resulting solution $u_{\gamma}^{f}(x)$ of \Eqref{eq:forward} determines the induced boundary current
\begin{equation}\label{eq:current}
g(x) = \gamma(x)\frac{\partial u_{\gamma}^{f}}{\partial n}(x) \quad \text{on } \partial\Omega\,,
\end{equation}
where $\frac{\partial}{\partial n} $ denotes the outward normal derivative. 
The mapping $f \mapsto g$ in \Eqref{eq:current} defines the Dirichlet-to-Neumann operator
\begin{equation}\label{eq:DtN_cont}
\Lambda_\gamma: H^{1/2}(\partial\Omega) \to H^{-1/2}(\partial\Omega)\,, 
\end{equation}
which takes a boundary voltage function $f$ on $\partial\Omega$ and returns the corresponding boundary current $g$. Here $H^{\pm 1/2}(\partial\Omega)$ is the standard Sobolev Hilbert spaces on $\partial\Omega$ with inner products
$\langle\cdot,\cdot\rangle_{H^{\pm 1/2}}$ and norms $\|\cdot\|_{H^{\pm 1/2}}$.

The weak form of \Eqref{eq:forward} can be formulated as find $u \in H^{1}(\Omega)$, $u|_{\partial \Omega} = f$ such that
\begin{equation}\label{eqn:weak_form}
    \int_{\Omega} \gamma(x) \nabla u \cdot \nabla v\diff x = 0,\quad \forall v \in H^{1}_{0}(\Omega)\,.
\end{equation}
Thus we could define DtN operator
\begin{equation}\label{weakform}
    \langle\Lambda_\gamma f, \psi\rangle = \int_{\partial\Omega} \gamma \frac{\partial u_{\gamma}^{f}}{\partial n} \psi \diff s\,.
\end{equation}
for $\psi \in H^{1/2}_0(\partial\Omega)$.



In EIT, we aim to recover $\gamma(x)$ from multiple data pairs $\{(f_i,g_i)\}$, where the input boundary voltage $f_i$ depends on experimental design. Ideally, if one could measure $g$ for all admissible $f$, $\gamma(x)$ can be recovered in theory \citep{uhlmann2009electrical,isakov2006inverse, astala2006calderon}. However, as infinitely many measurements are not attainable, using a finite number of measurement pairs $(f_i,g_i)$ is of more interest in practice. 


\paragraph{Finite element discretizations} 
We approximate $\Omega$ by a polygonal domain $\Omega_h$ with a shape-regular triangulation $\mathcal T_h$ and use the linear ($P_1$) finite element space $V_h\subset H^1(\Omega_h)$ of continuous, piecewise-linear functions on $\mathcal T_h$. Let $\{\phi_k\}_{k=1}^{N_{\text{el}}}$ be the nodal basis of $V_h$, and split the indices into interior nodes $\mathcal{I}$ with $N_{\mathcal{I}}=|\mathcal{I}|$ and boundary nodes $\mathcal{B}$ with $N_{\mathcal{B}}=|\mathcal{B}|$. Define the subspace $V_h^0 = \{v^h \in V_h: v^h |_{\partial\Omega_h}=0 \}$ and the trace space $V_h(\partial \Omega_h) = \operatorname{span}\{\phi_j |_{\partial\Omega_h}, j \in \mathcal{B}\}$.
Given boundary data $f^h \in V_h(\partial \Omega_h)$, \Eqref{eqn:weak_form} now can be discretized as follows:
\[
  \text{Find } u^h \in V_h \text{ with } u^h|_{\partial \Omega_h}=f^h,  \text{ such that }a(u^h, v^h) = \int_{\Omega_h} \gamma \nabla u^h \cdot \nabla v^h \diff x=0, \quad \forall v^h \in V_h^0\,.
\]
We assemble the stiffness matrix $\mathbf{K} \in \mathbb{R}^{N_{\text{el}} \times N_{\text{el}}}$ with entries $\mathbf{K}_{ij}=a(\phi_i, \phi_j)$ and denote by $\mathbf{u}$, $\mathbf{f}$ and $\mathbf{g}$ the corresponding vectors by tabulating its coefficient of the FEM basis $\{\phi_i\}_{i \in \mathcal{B} \cup \mathcal{I}}$ for $u^h, f^h$ and $g^h$ respectively.

We partition $\mathbf{K}$ and $\mathbf{u}$ with index sets $\mathcal{I}$ and $\mathcal{B}$ with
\begin{equation*}
\mathbf{K}=
\begin{bmatrix}
    \mathbf{K}_{\mathcal{I}\mathcal{I}} & \mathbf{K}_{\mathcal{I}\mathcal{B}} \\
    \mathbf{K}_{\mathcal{BI}} & \mathbf{K}_{\mathcal{BB}}
\end{bmatrix} \quad \text{and} \quad \qquad\mathbf{u} = \begin{bmatrix}
    \mathbf{u_{\mathcal{I}}}\\ \mathbf{u}_{\mathcal{B}}
\end{bmatrix} \,.
\end{equation*}

Let the discrete Neumann data vector $\mathbf{g} \in \mathbb{R}^{N_{\mathcal{B}}}$ be defined by $g_j = a(u^h, \phi_j)$ for $j = 1, \dots, N_\mathcal{B}$. Eliminating the interior degrees of freedom yields the block relation
\begin{equation}
\label{eq:g_lambda}
\mathbf{g}=\mathbf{K}_{\mathcal{B B}}\mathbf{f}+ \mathbf{K}_{\mathcal{BI}} \mathbf{u}_{\mathcal{I}}=\underbrace{\left(\mathbf{K}_{\mathcal{B B}}-\mathbf{K}_{\mathcal{BI}} \mathbf{K}_{\mathcal{II}}^{-1} \mathbf{K}_{\mathcal{I B}}\right)}_{\mathbf{\Lambda}_\gamma} \mathbf{f} \,.  
\end{equation}
Therefore, the discrete DtN matrix $\mathbf{\Lambda}_{\gamma}$ is the Schur complement of the stiffness matrix $\mathbf{K}$ with respect to the interior block $\mathbf{K}_{\mathcal{II}}$. If the $k$-th experiment impose the canonical boundary voltage $\mathbf{f} = \mathbf{e}_k$, then $\mathbf{g} = \mathbf{\Lambda}_\gamma \mathbf{e}_k$, so the corresponding Neumann measurement vector is exactly the $k$-th column of $\mathbf{\Lambda}_\gamma$. 
Consequently, under the canonical set of boundary excitations, $\mathbf{\Lambda}_\gamma$ coincides with the complete discrete current-measurement table on the prescribed boundary node set. 
More generally, for a multi-source experiment, we collect all boundary voltage vectors $\{\mathbf{f}_k\}$ as the columns of a matrix $\mathbf{F} \in \mathbb{R}^{N_\mathcal{B}\times N_\mathcal{B}}$, and the associated Neumann measurements are then assembled as $\mathbf{\Lambda}_\gamma \mathbf{F}$. Without loss of generality, we restrict attention to the canonical experiment design $\mathbf{F} = \mathbf{I}$, in which case the measurement matrix coincides with the DtN matrix.

\paragraph{Partial observations and data completion}
 Ideally, if all chosen boundary nodes are applied and the currents are measured at all boundary nodes, one would fully determine the $N_{\mathcal{B}} \times N_{\mathcal{B}}$ DtN matrix $\mathbf{\Lambda}_{\gamma}$. However, practical and economic constraints often prevent running all possible experiments or taking all desired measurements. In practice, only a small number of experiments can be
conducted, and in each experiment, only a small number of measurements can be
taken, which results in only a subset of DtN matrix entries being measured. Let $\mathbf{M}_s\in\{0,1\}^{N_{\mathcal B}\times N_{\mathcal B}}$ be a binary mask matrix with $(\mathbf{M}_s)_{ij}=1$ if the entry $(i,j)$ is observed and $0$ otherwise, where $s=\frac{\| \mathbf{M}_s\|_0} {N^2_\mathcal{B}}$ is the sampling rate, i.e., the portion of observed entries over all entries. Thus the \emph{configuration} of how we obtain the DtN measurements is encoded by $\mathbf{M}_s$, and the observed (incomplete) DtN measurements are denoted by
\[
\mathbf{\Lambda}^{\text{o}}_\gamma \;=\; \mathbf{M}_s \odot \mathbf{\Lambda}_\gamma\,,
\]
where $\odot$ denotes the Hadamard product. We emphasize that \(\mathbf M_s\) is an abstract encoding of the acquisition pattern and may represent highly structured missingness (e.g., missing columns corresponding to unperformed excitations, or subsampled receiver sets), not necessarily i.i.d.\ entrywise sampling.


Given the masked DtN matrix $\mathbf{\Lambda}^{\mathrm{o}}_\gamma$ and mask $ \mathbf{M}_s$, our completion task is to infer the missing entries, i.e. estimate a completed matrix $\widehat{\mathbf{\Lambda}}_\gamma$ that is close to the true $\mathbf{\Lambda}_\gamma$. 



\section{Method}\label{sec:method}

In this section, we will formulate the data completion problem by conditional sampling solved by diffusion model. As multiple completions may plausibly explain the same partial observations, we model $\mathbf{\Lambda}_\gamma$ as a random matrix whose randomness comes from some distribution of $\gamma$ and the goal is to generate samples from the learned conditional distribution $p\!\left(\mathbf{\Lambda}_{\gamma}\,\middle|\,\mathbf{\Lambda}^{\mathrm{o}}_{\gamma}, \mathbf{M}_s\right)$. 
The generated posterior samples can then be used to estimate the full DtN matrix $\hat{\mathbf{\Lambda}}_\gamma$, and standard inverse solvers including iterative solvers, optimization-based reconstruction, or deep learning-based methods can be applied to recover conductivity $\gamma$.
The key is to ensure that $\hat{\mathbf{\Lambda}}_\gamma$ is as close as possible to the true underlying operator $\mathbf{\Lambda}_\gamma$, thereby achieving better results in the downstream inversion than the same solver without completion.

\subsection{Diffusion Model} 

Given i.i.d.\ samples from a target distribution $p_{\text{target}}$, a diffusion model seeks to generate new samples whose distribution approximates $p_{\text{target}}$.
Let $\{\alpha_t\}_{t=1}^T\subset(0,1]$ and define $\bar{\alpha}_t\coloneqq \prod_{i=1}^t\alpha_i$.
Define the forward Markov chain by \(X_0\sim p_{\text{target}}\) and, for \(t=1,\dots,T\), \(X_t=\sqrt{\alpha_t}\,X_{t-1}+\sqrt{1-\alpha_t}\,W_t\) with \(W_t\stackrel{\text{i.i.d.}}{\sim}\mathcal N(0,I_d)\); equivalently, \(X_t\mid X_0\sim \mathcal N(\sqrt{\bar\alpha_t}\,X_0,(1-\bar\alpha_t)I_d)\), hence \(X_t=\sqrt{\bar\alpha_t}\,X_0+\sqrt{1-\bar\alpha_t}\,\overline W\) for \(\overline W\sim\mathcal N(0,I_d)\).
This construction progressively perturbs $X_0$ with Gaussian noise over $t=1,\ldots,T$.
The noise schedule $\{\alpha_t\}$ is chosen so that $\bar{\alpha}_T\approx 0$, in which case the marginal distribution of $X_T$ is close to a standard Gaussian, i.e., $p_{X_T}\approx \mathcal N(0,I_d)$.


The forward process is reversible as shown in classical results from stochastic differential equation (SDE) \citep{anderson1982reverse}, which inspires the sampling algorithm from diffusion model. 
The key aspect of the sampling algorithm from diffusion model is to construct a reverse process $Y_t$, which starts from pure Gaussian noise and can gradually convert to a new sample $Y_0$ such that $p_{Y_0} \approx p_{\text{target}}$. In particular, the reverse process relies on the availability of the so called \emph{score function} $s_t^\star(\cdot)\coloneqq \nabla\log p_{X_t}(\cdot)$, thus diffusion model is also called score-based generative model \cite{song2020score}. A classical choice of reverse process is the Denoising Diffusion Probabilistic Model (DDPM) \citep{ho2020denoising} sampler as follows, which can be treated as a discretization of reverse SDE. 
\begin{definition}[Unconditional DDPM sampler]\label{def:ddpm}
Given score estimator $s_t(\cdot)$ that approximate the true score function $s_t^{\star}(\cdot)$, i.e., $s_t(\cdot) \approx s_t^{\star}(\cdot)$ and the same noise schedule $\{\alpha_t\}$ as in forward process, DDPM sampler is defined as
\begin{equation}
Y_T\sim\mathcal N(0,I_d),\qquad
Y_{t-1}
=\frac{1}{\sqrt{\alpha_t}}\Big(Y_t+(1-\alpha_t)\,s_t(Y_t)+\sigma_t W_t\Big),
\quad W_t\stackrel{\text{i.i.d.}}{\sim}\mathcal N(0,I_d)\,,
\end{equation}
with
\begin{equation}
\sigma_t
\;=\;\sqrt{\frac{(\alpha_t-\bar{\alpha}_t)(1-\alpha_t)}{1-\bar{\alpha}_t}}\,,\qquad t=T,\ldots,1\,.
\end{equation}
\end{definition}
When a score estimator $s_t(\cdot) \approx s_t^\star(\cdot)$ is available, DDPM sampler \ref{def:ddpm} provides a way to generate new samples from the target distribution \citep{ho2020denoising}. In practice, the score estimator is parametrized by a neural network $s_{\theta}(\cdot, t)$, and the parameters $\theta$ can be learned through the \emph{score matching loss} \citep{ho2020denoising, song2020score, vincent2011connection}:
\begin{equation}
\label{eq:score-matching} 
    \min_{\theta} \mathbb{E}_{t \sim \mathrm{Unif}{(1,\dots,T)}, X_0 \sim p_{\text{target}}, X_t \sim p_{X_t \mid X_0}} \lambda(t)\| s_{\theta}(X_t,t) - \nabla_{x_t} \log p(X_t \mid X_0) \|^2\,,
\end{equation}
where $t$ follows a uniform distribution over $1,2,\dots,T$, $\lambda(t)$ is a positive weighting function, and $\nabla\log p(X_t \mid X_0)$ can be computed \emph{analytically} from the forward diffusion process. Concretely, the expectation can be approximated by the Monte Carlo estimation. Each training iteration (1) draws a batch $\{ X_0^{(i)}\}_{i=1}^B$ with size $B$ from dataset $\mathcal{D}_{\text{train}}$, (2) samples an independent time step $t_i \sim \mathcal{DU}\{1,\dots,T\}$ and noise $W^i \sim \mathcal{N}(0,I_d)$ for each sample, and (3) forms a noised input $X_{t_i}^{(i)} = \sqrt{\bar\alpha_{t_i}}\,X_0^{(i)} + \sqrt{1-\bar\alpha_{t_i}}\,W^{i}$. We then compute the regression target $s_{t_i}(X_{t_i}^{(i)}) = -\dfrac{X_{t_i}^{(i)}-\sqrt{\bar\alpha_{t_i}}\,X_0^{(i)}}{1-\bar\alpha_{t_i}}$, evaluate the weighted mean-squared error $\lambda(t_i)\| s_{\theta}(X^{(i)}_{t_i},t_i) - s_{t_i}(X_{t_i}^{(i)}) \|^2$ over the batch, and update $\theta$ with standard deep learning solver, e.g. AdamW \citep{loshchilov2017decoupled}. 
At the population level, \Eqref{eq:score-matching} is minimized when $s_\theta(\cdot,t)$ matches the true marginal score $s_t^\star(x)\coloneqq\nabla_x\log p_t(x)$ (where $p_t$ is the law of $X_t$ induced by the forward diffusion); in well-specified and identifiable parametric settings, minimizers of the empirical (Monte Carlo) counterpart converge in probability to a population minimizer and hence yield a consistent score estimator \citep{hyvarinen2005estimation,song2020score}.

With standard assumptions on sufficiency of the model capacity, the optimal solution of ~\Eqref{eq:score-matching}, denoted as $\theta^\star$, is a consistent estimator, i.e., it converges in probability towards
the true value of $\theta$ when sample size goes infinity \citep{song2020score, hyvarinen2005estimation}.

\subsection{Conditional Diffusion Models for Data Completion}
\label{sec:proposed_method}
Diffusion model is also flexible in incorporating a conditional covariate, denoted as $Z$, to guide the generation of new samples.  For each fixed $Z=z$, DDPM sampler~\ref{def:ddpm} can be extended to conditional generation when the conditional score estimator $s_t(\cdot, z) \approx s_t^{\star}(\cdot,z)$ is given, where $s_t^{\star}(\cdot,z)=\nabla \log p_{X_t|z}(\cdot)$. The following conditional DDPM sampler can generate a new sample $Y_0$ such that $p_{Y_0} \approx p_{X_0 \mid Z=z}$.  
\begin{definition}[Conditional DDPM sampler]\label{def:cond-DDPM-sampler}
For a fixed $Z=z$, given conditional score estimator $s_t(\cdot,z)$ that approximate the true conditional score function $s_t^{\star}(\cdot,z)$, i.e. $s_t(\cdot,z) \approx s_t^{\star}(\cdot,z)$, the conditional DDPM sampler is defined as 
\begin{equation}
Y_T\sim\mathcal N(0,I_{d}),\qquad
Y_{t-1}=\frac{1}{\sqrt{\alpha_t}}\Big(Y_t+(1-\alpha_t)\,s_t(Y_t, z)+\sigma_tW_t\Big),
\ \ W_t\stackrel{\text{i.i.d.}}{\sim}\mathcal N(0,I_{d})\,.
\end{equation}

Here $\alpha_t$, $\bar{\alpha}_t$, $\sigma_t$ are the same as defined in unconditional DDPM sampler~\ref{def:ddpm}.
\end{definition}

The conditional score function can also be parametrized by a neural network $s_{\theta}(\cdot, z, t)$, and learned by following conditional score matching,
\begin{equation}
\label{eq:cond-score-matching}
    \min_{\theta} \mathbb{E}_{t \sim \mathrm{Unif}{(1,\dots,T)}, (X_0, Z) \sim p_{(X_0,Z)}, X_t \sim p_{X_t \mid X_0}} \lambda(t)\| s_{\theta}(X_t,Z,t) - \nabla_{x_t} \log p(X_t \mid X_0) \|^2\,,
\end{equation}
where $p_{(X_0, Z)}$ is the joint distribution of target $X_0$ and condition $Z$. 
Although $p_{X_t \mid Z}$ does not appear in the loss function explicitly, it has been proved in \citet{batzolis2021conditional,vincent2011connection} that $s_{\theta^\star}(X, Z, t)$ is a consistent estimator of $\nabla\log p(X_t \mid Z)$ under some standard assumptions, where $\theta^\star$ is the minimizer of ~\Eqref{eq:cond-score-matching}. In practice we approximate the expectation in \Eqref{eq:cond-score-matching} with the same Monte Carlo scheme as in the unconditional case, except that each training example carries its condition $Z$ and the network receives $(X_t,Z,t)$ as input. 

Recall that in EIT data completion, our goal is to sample $p(\mathbf{\Lambda}_\gamma \mid \mathbf{\Lambda}_{\gamma}^{\text{o}},\mathbf{M})$ for any given partial DtN measurements $\mathbf{\Lambda}_{\gamma}^{\text{o}}$ and mask $\mathbf{M}$. We instantiate the conditional diffusion model by setting target variable $X_0=\mathbf{\Lambda}_\gamma$ and conditional variable $Z=(\mathbf{\Lambda}_{\gamma}^{\text{o}},\mathbf{M})$. Throughout, we identify a DtN matrix $X_0\in\R^{N_{\mathcal B}\times N_{\mathcal B}}$ with its vectorization $\mathrm{vec}(X_0)\in\R^{d}$, $d=N_{\mathcal B}^2$; the forward/noising and reverse/sampling updates are applied to $\mathrm{vec}(X_0)$, and the matrix form is only a reshape used for notation and implementation.
Concretely, suppose various configurations $\mathbf{M}$ follow some distribution $\mathcal{M}$, we draw full DtN measurements $\{\mathbf{\Lambda}_{\gamma_i}\}_{i=1}^B$ from $\mathcal{D_{\text{train}}}$ and $\{\mathbf{M}^{(i)}\}_{i=1}^B \sim\mathcal{M}$. For each $i$, we form the pair sample 
\[
\{X_0^{(i)}, Z^i\} = \{\mathbf{\Lambda}_{\gamma_i},(\mathbf{M}^{(i)}\odot \mathbf{\Lambda}_{\gamma_i}, \mathbf{M}^{(i)})\}\,,
\] which yields samples from the joint distribution $p(X_0, Z)$. The conditional score $s_\theta(X_t,Z,t)$ is then trained exactly as in \Eqref{eq:cond-score-matching}. At test time, for a given partially observed DtN measurements $\mathbf\Lambda^{\mathrm{o}}_{\gamma}=\mathbf M^{\text{test}}\odot\mathbf\Lambda_{\gamma}^{\text{test}}$, we set $Z_{\text{test}}=(\mathbf\Lambda^{\mathrm{o}}_{\gamma}, \mathbf{M}^{\text{test}})$ and run the conditional DDPM sampler \ref{def:cond-DDPM-sampler} to generate $\widehat{\mathbf{\Lambda}}_\gamma = Y_0$ such that $\widehat{\mathbf{\Lambda}}_\gamma \approx \mathbf\Lambda_{\gamma}^{\text{test}}$. The training and sampling pseudo-codes are summarized in Algorithm \ref{alg:train-conditional-score} and Algorithm \ref{alg:complete-then-reconstruct}, respectively. 

\begin{algorithm}[H]
\caption{Training the conditional score network for DtN measurements completion}
\label{alg:train-conditional-score}
\begin{algorithmic}[1]
\Require Fully observed DtN dataset $\mathcal{D_{\text{train}}}=\{\mathbf{\Lambda}_{\gamma_i}\}_{i=1}^{N_{\text{train}}}$; measurement configuration distribution $\mathcal{M}$; schedule $\{\alpha_t\}_{t=1}^T$, weight $\lambda(t)$, batch size $B$; conditional score network $s_\theta(X_t,Z,t)$ to be optimized and optimizer
\While{not converged}
  \State Sample $\{\mathbf{\Lambda}_{\gamma_i}\}_{i=1}^B$ from $\mathcal{D}_{\text{train}}$, and $\{\mathbf{M}^{(i)}\}_{i=1}^B \sim \mathcal{M}$
  \State Set $X_0^{(i)} \gets\mathbf{\Lambda}_{\gamma_i}$ and conditions $Z^{i} \gets (\mathbf{M}^{(i)} \odot \mathbf{\Lambda}_{\gamma_i}, \mathbf{M}^{(i)})$
  \State Sample batch $\{t_i\}_{i=1}^B \sim \mathcal{DU}\{1,\dots,T\}$ and noise $\{W^{i}\}_{i=1}^B \sim \mathcal{N}(0,I_d)$ with same shape as $\mathbf{\Lambda}_{\gamma_i}$
  \State $X_{t_i}^{(i)} \gets \sqrt{\bar\alpha_{t_i}}\,X_0^{(i)} + \sqrt{1-\bar\alpha_{t_i}}\,W^{i}$ \Comment{forward diffusion}
  \State $s_{t_i}^\star(X_{t_i}^{(i)}) \gets -\dfrac{X_{t_i}^{(i)}-\sqrt{\bar\alpha_{t_i}}\,X_0^{(i)}}{1-\bar\alpha_{t_i}}$ \Comment{analytic score of $p(X^{(i)}_{t_i}\mid X_0^{(i)})$}
  \State $\mathcal{L} \gets \dfrac{1}{B}\sum_{i=1}^B \lambda(t_i)\,\big\|s_\theta(X_{t_i}^{(i)},Z^{i},t_i)-s_{t_i}^\star(X_{t_i}^{(i)})\big\|^2$
  \State Update $\theta$ using optimizer step to minimize $\mathcal{L}$
\EndWhile
\State \textbf{return} $\theta$ \Comment{trained conditional score network}
\end{algorithmic}
\end{algorithm}

\begin{algorithm}[H]
\caption{DtN measurements completion via conditional diffusion model}
\label{alg:complete-then-reconstruct}
\begin{algorithmic}[1]
\Require Observed (masked) DtN measurements $\mathbf{\Lambda}^{\text{o}}_{\gamma}$ and its mask $\mathbf{M}^{\text{o}}$; trained $s_\theta$; schedule $\{\alpha_t\}_{t=1}^T$; conditional DDPM solver \ref{def:cond-DDPM-sampler}; 

\State Initialize $Y_T \sim \mathcal{N}(0,I_d)$ with same shape as $\mathbf{\Lambda}_\gamma^{\text{o}}$, and set the condition $Z \gets (\mathbf{\Lambda}_\gamma^{\text{o}},\mathbf{M}^{\text{o}}) $
\For{$t=T,\,T\!-\!1,\,\dots,\,1$}
  \State Sample $W_t \sim \mathcal{N}(0,I)$ with same shape as $\mathbf{\Lambda}_\gamma^{\text{o}}$
  \State $Y_{t-1} \gets \dfrac{1}{\sqrt{\alpha_t}}\Big(Y_t + (1-\alpha_t)\, s_\theta(Y_t, Z, t) + \sigma_t W_t\Big)$ \Comment{Conditional DDPM \ref{def:cond-DDPM-sampler} update}
\EndFor
\State $\widehat{\mathbf{\Lambda}}_\gamma \gets Y_0$

\State \textbf{return} $\widehat{\mathbf{\Lambda}}_\gamma $ 
\end{algorithmic}
\end{algorithm}

\subsection{Solve the inverse problem by data completion}
\begin{figure}[!htb]
    \centering
    \includegraphics[width=1\linewidth]{diagram.png}
\caption{Three strategies of solving EIT. Direct reconstruction \textbf{(first row)} from sparse measurements yields degraded reconstruction. We propose to first complete the DtN measurements and then reconstruct \textbf{(second row)}, yielding image quality close to those from a direct reconstruction from full measurements \textbf{(third row)}.
}
\label{fig:diagram}
\end{figure}
We illustrate our framework for solving the EIT inverse problem in Figure \ref{fig:diagram}. The figure compares three reconstruction strategies: direct reconstruction from sparse DtN measurements (first row), reconstruction from DtN measurements that are first completed and then passed to the same solver (second row), and direct reconstruction from fully sampled DtN measurements with the same solver (third row). The completed DtN measurements produced by our method are designed to be plug-and-play. Once the partially observed operator is completed to $\widehat{\mathbf{\Lambda}}_\gamma$, it can be passed to any inverse solver that presupposes full boundary data without modification. In our experiments, rather than using a classical PDE-based inverse solver, we adopt deep learning-based inverse solvers for their stable and excellent performance. Moreover, these deep learning models offer a unified pipeline to handle both full and partial measurements, making it particularly transparent to isolate and evaluate the effect of the data-completion step.




\section{Error Analysis}\label{sec:theory}
In this section, we provide an end-to-end error analysis for our algorithm. We are able to quantify the total variation (TV) discrepancy between the completed DtN measurements from our algorithm and that of the ground truth measurements. 
This section is organized as follows. In Section~\ref{sec:K-and-param} we first introduce the admissible class \(K\) of conductivities of $n_{\mathrm{v}}$-polygon shape, and its vertex parametrization $v$. 
We next show $K$ is compact and as a consequence we can construct finitely many local charts that cover \(K\) and yield bounded parameter patches $E_i$. 
In Section~\ref{sec:chartwise_DtN}, we show that the forward map \(v\mapsto\Lambda_{\gamma(v)}\) is \(C^1\) in operator norm in each patch. In Section~\ref{sec:chartwise_lipschitz_discretized_DtN}, we define a finite-dimensional DtN measurement map $F_N$ and establish its patchwise Lipschitz-type stability. In Section~\ref{sec:entropy-bdd-FN}, we leverage these patchwise stability bounds to control covering numbers for the induced measurement set. We then apply a conditional DDPM guarantee in Section~\ref{sec:ddpm_FNK} to translate these geometric controls together with a score-error assumption into the non-asymptotic DtN completion bound of Theorem~\ref{thm:cond-ddpm}. Overall, this section yields an end-to-end, non-asymptotic TV error guarantee for DtN completion over the polygonal conductivity class.

\paragraph{Notations.}
We denote by 
$\dist(\cdot,\cdot)$ the Euclidean distance between points or subsets in $\mathbb{R}^2$, and $|\cdot|$ the Euclidean norm in $\mathbb{R}^2$. 
We also use $|E|$ to denote the Lebesgue measure of a measurable set $E$ when no ambiguity arises.
For nonempty sets $A,B\subset\R^2$, the Hausdorff distance is  
\[
d_H(A,B)\coloneqq \max\Big\{\sup_{a\in A} \inf_{b \in B}\dist(a,b),\ \sup_{b\in B} \inf_{a \in A} \dist(b,a)\Big\}.
\]
For measurable sets $E,F\subset\R^2$, the symmetric difference is
$E\Delta F\coloneqq (E\setminus F)\cup(F\setminus E)$. We denote by $B_2(x,r)\coloneqq \{y\in\R^2:\ |y-x|\le r\}$ the closed Euclidean ball, and by
$\diam(\Omega)\coloneqq \sup\{|x-y|:\ x,y\in\Omega\}$ the diameter of a set $\Omega$.
For a simple polygon $P(v)$ with $n_v\in \mathbb{N}$ vertices $v=(v_i)_{i=0}^{n_{\mathrm v}-1}$, we write its perimeter as
\[
\Per(P(v))\coloneqq \sum_{i=0}^{n_{\mathrm v}-1}|v_{i+1}-v_i|,\qquad (v_{n_{\mathrm v}}\coloneqq  v_0)\,.
\]
Throughout the paper, we always assume the vertices $v_i$'s for a polygon $P(v)$ are ordered in the counterclockwise order.
For a set $E\subset\R^{2n_{\mathrm v}}$, we write $\conv(E)$ for its convex hull.
Finally, for a measurable set $E\subset\Omega$, $\chi_{E}$ denotes its indicator function.

\subsection{Admissible conductivities and parametrization}\label{sec:K-and-param}

Consider any integer $n_{\mathrm v}\ge 3$,
the vertex list $v=(v_0,\dots,v_{n_{\mathrm v}-1})\in(\R^2)^{n_{\mathrm v}}$ consists pairwise distinct points with $v_{n_{\mathrm v}}:=v_0$. We define the closed polygonal chain
\[
\partial P(v)\coloneqq \bigcup_{i=0}^{n_{\mathrm v}-1}[v_i,v_{i+1}],
\qquad
[v_i,v_{i+1}]\coloneqq \{(1-t)v_i+t v_{i+1}:t\in[0,1]\}.
\]
We denote by $\mathcal V^{n_{\mathrm v}}$ the set of vertex lists $v$ such that $\partial P(v)$ is a
simple (non self-intersecting) closed polygonal curve and the vertices are listed counterclockwise.
For $v\in\mathcal V^{n_{\mathrm v}}$, the Jordan curve theorem yields a unique bounded connected component
$\operatorname{int}(P(v))$ of $\R^2\setminus \partial P(v)$; we define the (closed) $n_{\mathrm v}$-gon $P(v)\coloneqq \overline{\operatorname{int}(P(v))}$. Note that the parametrization $v\mapsto P(v)$ is not injective due to the ordering of vertices (cyclic shifts and reversal).

For $v\in\mathcal V^{n_{\mathrm v}}$, we denote by $\beta_i(v)$
the interior angle of $P(v)$ at the vertex $v_i$. We also define the polygon norm
\[
\|v\|_{\mathrm{poly}}\coloneqq \max_{0\le i\le n_{\mathrm v}-1}|v_i|,
\qquad
\|v-w\|_{\mathrm{poly}}\coloneqq \max_{0\le i\le n_{\mathrm v}-1}|v_i-w_i|.
\]
Identifying $v$ with an element of $\R^{2n_{\mathrm v}}$ by concatenation, we have
$\|v\|_{\mathrm{poly}}\le \|v\|_2 \le \sqrt{n_{\mathrm v}}\,\|v\|_{\mathrm{poly}}$.

\begin{assumption}[Admissible polygon classes]\label{ass:admissible_poly}
Fix integers $n_{\mathrm v}\ge 3$ and constants $\beta_0\in(0,\pi/2)$, $d_0>0$, $d_1>0$.
Let $\Omega\subset\R^2$ be the unit disk.
\begin{definition}[Admissible polygon class]\label{def:admiss}
Define the \emph{admissible polygon class} $\mathcal A$ as the set of all polygons $P=P(v)$ with their vertices $v\in\mathcal V^{n_{\mathrm v}}$
such that:
\begin{enumerate}
\item[(A1)] $P\subset\Omega$ and $\dist(P,\partial\Omega)\ge d_0$\,;
\item[(A2)] every side length satisfies $|v_{i+1}-v_i|\ge d_1$ for all $i$ (indices mod $n_{\mathrm v}$)\,;
\item[(A3)] every interior angle satisfies $\beta_0\le \beta_i(v)\le \pi-\beta_0$ for all $i$; 
\end{enumerate}
\end{definition}
\begin{remark}
    The above assumptions in ~Definition \ref{def:admiss} imply that $\mathcal{A}$ consists of convex polygons with nonzero side length that are away from the boundary $\partial\Omega$.
    These assumptions satisfy the geometric hypotheses used in \citealp[Prop.~3.3]{beretta2022global}. In particular, since we only consider on convex polygon, the uniform Lipschitz boundary assumption~\ref{def: lipschitz_bdy_class} in \citealp[Prop.~3.3]{beretta2022global} can be derived by A1--A3. We show this in the Appendix~\ref{sec:lipschitz_bdy_class}. 
\end{remark}
\begin{definition}[Relaxed admissible polygon classes]
Define the \emph{relaxed class} $\mathcal A_{1/2}$ by requiring the same conditions as above,
but with \emph{strict} inequalities, and the parameters are relaxed to half of their values in $\mathcal{A}$:
\begin{enumerate}
\item[(A1$'$)] $P\subset\Omega$ and $\dist(P,\partial\Omega) > d_0/2$;
\item[(A2$'$)] $|v_{i+1}-v_i| > d_1/2$ for all $i$ (indices mod $n_{\mathrm v}$)\,;
\item[(A3$'$)] every interior angle satisfies $\frac{\beta_0}{2}< \beta_i(v)< \pi-\frac{\beta_0}{2}$ for all $i$;
\end{enumerate}
\end{definition}
\end{assumption}

We say $v$ is admissible if $P(v) \in \mathcal{A}$ and  relaxed admissible if $P(v) \in \mathcal{A}_{1/2}$. 
For an admissible polygon $P$ we consider conductivity
\[
\gamma_{P}\coloneqq 1+(\kappa-1)\chi_{P}\,.
\]

Finally, we define the relaxed and admissible conductivity classes
\[
M\coloneqq \{\gamma_P:\ P\in\mathcal A_{1/2}\},
\qquad
K\coloneqq \{\gamma_P:\ P\in\mathcal A\}\subset M\,.
\]

\begin{remark}[A priori data]\label{rem:apriori_poly}
All constants introduced below depend only on the \emph{a priori data}
\[
(\Omega,\kappa,n_{\mathrm v},d_0,d_1, \beta_0),
\]
and are independent of the discretization level $N$ and the number of diffusion steps $T$ in later Section \ref{sec:chartwise_lipschitz_discretized_DtN}, Section \ref{sec:entropy-bdd-FN} and Section \ref{sec:ddpm_FNK}.
\end{remark}


\begin{lemma}[Local admissibility radius]\label{lem:rho_s}
Fix $P^\ast\in\mathcal A_{1/2}$ and let $v^\ast$ be its counterclockwise ordered vertex list.
Then there exists $\rho_{P^\ast}>0$ such that whenever $\|v-v^\ast\|_{\mathrm{poly}}<\rho_{P^\ast}$,
the polygonal chain $\partial P(v)$ is simple and $P(v)\in\mathcal A_{1/2}$.
\end{lemma}
Lemma~\ref{lem:rho_s} shows for any relaxed admissible polygon $P^\ast$, there is a small ball in the vertex space such that all polygons with vertices in the ball stays
in the relaxed class. We prove Lemma~\ref{lem:rho_s} in Appendix~\ref{sec:proof_rho_s}. The following Lemma~\ref{lem:chart-M} constructs the atlas for $M$, which follows the 8.1.3-8.1.5 in \citep{alberti2022inverse} and is proved in Appendix~\ref{sec:chart-M}.

\begin{lemma}\label{lem:chart-M}
Fix $P^\star\in\mathcal A_{1/2}$ with counterclockwise ordered vertices
$v^\star=(v_i^\star)_{i=0}^{n_{\mathrm v}-1}$, and let $\rho_{P^\star}>0$ be as in
Lemma~\ref{lem:rho_s}. Define
\[
R_{P^\star}\coloneqq \frac12\min\Big\{\frac14\min_{i\neq j}|v_i^\star-v_j^\star|,\ 
\dist(P^\star,\partial\Omega)-\frac12 d_0,\ \rho_{P^\star}\Big\}>0,
\]
\[
\mathcal O_{P^\star}\coloneqq \Big\{v\in\R^{2n_{\mathrm v}}:\ \|v-v^\star\|_{\mathrm{poly}}<R_{P^\star}\Big\},
\qquad
U_{P^\star}\coloneqq \Big\{\gamma_{P(v)}:\ v\in\mathcal O_{P^\star}\Big\}, 
\]
and
\[
\varphi_{P^\star}:U_{P^\star}\to\mathcal O_{P^\star},\qquad
\varphi_{P^\star}(\gamma_{P(v)})\coloneqq v .
\]
Then
\begin{enumerate}
\item[(i)] $U_{P^\star}$ is open in $M$ with respect to the subspace topology induced by $L^1(\Omega)$.
\item[(ii)] $\varphi_{P^\star}$ is well-defined and is a homeomorphism from $U_{P^\star}$ onto
$\mathcal O_{P^\star}$. 
\end{enumerate}
\end{lemma}


In particular, the choice of $R_{P^\star}$ ensures (1) the neighborhood of $P^\star$ is still inside $\mathcal{A}_{1/2}$ and (2) $\varphi$ is well defined and injective by ruling out the nonuniqueness caused by reordering vertex. 


\begin{lemma}[Compactness of $K$]\label{lem:K_compact_poly}
The admissible conductivity set $K$ is compact in $L^1(\Omega)$.
\end{lemma}
We show the compactness of the admissible set in Appendix~\ref{sec:K_compact_poly}, which allows us to extract an finite cover of $K$ by leveraging the local parameter neighborhoods defined in Lemma~\ref{lem:chart-M}.

\begin{lemma}[Finite patch cover and compact parameter patches]\label{lem:finite-atlas}
For each $\gamma_P\in K$, let $v^P:=\varphi_P(\gamma_P)$ and define the shrunken neighborhood
\[
V_{P}:=\Big\{\gamma\in U_P:\ \|\varphi_P(\gamma)-v^{P}\|_{\mathrm{poly}}<\tfrac12 R_{P}\Big\}\subset U_P .
\]
Then there exist $P_1,\dots,P_m\in\mathcal A$ such that
\[
K\subset \bigcup_{i=1}^m V_{P_i}.
\]
For each $i$, define the parameter patch
\[
E_i:=\varphi_{P_i}\bigl(K\cap \overline{V}_{P_i}\bigr)\subset \R^{2n_{\mathrm v}},
\]
where $\overline{V}_{P_i}$ denotes the closure in $M$.
Then each $E_i$ is compact in $\R^{2n_{\mathrm v}}$. Moreover, if we set
\[
B_i:=\Big\{v\in\R^{2n_{\mathrm v}}:\ \|v-v^{P_i}\|_{\mathrm{poly}}\le \tfrac12 R_{P_i}\Big\},
\]
then $E_i\subset B_i\subset \varphi_{P_i}(U_{P_i})$, hence $\bigcup_{i=1}^m E_i$ is bounded and $\operatorname{conv}(E_i)\subset \varphi_{P_i}(U_{P_i})$ for each  $i=1,\dots,m$.

\end{lemma}

The finite many parameter patches $E_i$ will help us to quantify the complexity of DtN measurements through the local regularity of mapping from vertex to DtN measurements. 


\subsection{Chartwise regularity of the DtN map}\label{sec:chartwise_DtN}
To transfer low-dimensional structure from vertices to measurements, we need a quantitative regularity statement for the forward map. The regularity of mapping $v \to \Lambda_{\gamma(v)}$ has been investigated in \citet{alberti2022inverse, beretta2017differentiability, beretta2022global}. Within each chart, the DtN operator depends $C^1-$smoothly on the vertex parameters. 
Consider
\[
Y \coloneqq  \mathcal{L}\!\big(H^{1/2}(\partial\Omega),\,H^{-1/2}(\partial\Omega)\big),
\qquad
\|A\|_Y \coloneqq  \sup_{\|f\|_{H^{1/2}(\partial\Omega)}=1}\|Af\|_{H^{-1/2}(\partial\Omega)}\,,
\]
Lemma~\ref{lem:DtN-C1-chart} states the regularity statement of the forward map chartwisely. 


\begin{lemma}[$C^1$-regularity of the DtN map in local coordinates]\label{lem:DtN-C1-chart}
Fix $P^\star\in\mathcal A_{1/2}$ and let $(U_{P^\star},\varphi_{P^\star})$ be the chart from Lemma~\ref{lem:chart-M}. Set $\mathcal O_{P^\star}\coloneqq \varphi_{P^\star}(U_{P^\star})\subset\mathbb R^{2n_{\mathrm v}}$ and for $v \in \mathcal{O}_{P^\star}$ define
\[
\widetilde\Lambda_{P^\star}:\mathcal O_{P^\star}\to Y,\qquad
\widetilde\Lambda_{P^\star}(v)\coloneqq \Lambda_{\gamma_{P(v)}}\,.
\]
Then $\widetilde\Lambda_{P^\star}\in C^1(\mathcal O_{P^\star};Y)$.
\end{lemma}


\begin{proof}[Proof of Lemma ~\ref{lem:DtN-C1-chart}]
By Lemma~\ref{lem:chart-M}, for all $v\in\mathcal O_{P^\star}$ the polygon $P(v)$ stays in $\mathcal A_{1/2}$ and the
counterclockwise vertex labeling is fixed in the chart.
For polygonal inclusions satisfying the a priori geometric constraints, the map $v\mapsto\Lambda_{\gamma_{P(v)}}$ admits
a shape derivative with respect to vertex perturbations, and this derivative is continuous in the operator norm of
$Y$. This is established in the differentiability results for polygonal conductivities (see, e.g.,
\citealp[Lem.~4.4 and Cor.~4.5]{beretta2017differentiability} and the discussion in \citealp[Sec.~8.1]{alberti2022inverse}).
Therefore $\widetilde\Lambda_{P^\star}$ is continuously Fr\'{e}chet differentiable on $\mathcal O_{P^\star}$.
\end{proof}

Since $\widetilde\Lambda_{P_i}\in C^1(\mathcal O_{P_i};Y)$, the derivative
$D\widetilde\Lambda_{P_i}:\mathcal O_{
P_i}\to \mathcal L(\R^{2n_{\mathrm{v}}},Y)$ is continuous.
Moreover, $\operatorname{conv}(E_i)\subset \mathcal O_{P_i}$ and $\operatorname{conv}(E_i)$ is compact.
Hence we can denote
\begin{equation}
    M_i \coloneqq  \sup_{v\in \operatorname{conv}(E_i)}\big\|D\widetilde\Lambda_{P_i}(v)\big\|_{\mathrm{op}} \ <\ \infty,
\qquad
M_*\coloneqq \max_{1\le i\le m}M_i\,.
\end{equation}


We will later use $M_i$ and $M_*$ to show the chartwise Lipschitz for discretized DtN measurements. For simplicity, we denote $\widetilde\Lambda_{i}\coloneqq \widetilde\Lambda_{P_i}$.



\subsection{Chartwise Lipschitz for the DtN measurements}\label{sec:chartwise_lipschitz_discretized_DtN}

We next specify the finite-dimensional representation of DtN measurements. Starting from the DtN operator $\Lambda_\gamma\in Y$, we will (i) restrict it to finite-dimensional subspaces via the $R_N$ and obtain DtN matrix as in \Eqref{eqn:R_N}, 
and (ii) vectorize and normalize the resulting DtN matrix to obtain a vector in Euclidean space, which is aligned to the training object in diffusion model. 
Combining all these yields our discretized measurement map $F_N:K\to\R^{d_N}$ in \Eqref{eqn:F_N}. 

Let $V_N^1\subset H^{1/2}(\partial\Omega)$ and $V_N^2\subset H^{-1/2}(\partial\Omega)$ be finite-dimensional Hilbert subspaces with inherited inner products, and denote
\[
n_1\coloneqq \dim V_N^1,\qquad n_2\coloneqq \dim V_N^2,\qquad N\coloneqq \min\{n_1,n_2\},\qquad d_N\coloneqq n_1n_2,\qquad Y_N\coloneqq \mathcal{L}(V_N^1,V_N^2)\,.
\]

For notational simplicity, in the sequel we focus on the square case $n_1=n_2=N$, which is also consistent with our numerical setup and in this case $d_N=N^2$.
Let $P_N^1:H^{1/2}(\partial\Omega)\to V_N^1$ and $P_N^2:H^{-1/2}(\partial\Omega)\to V_N^2$
be the orthogonal projections and define 
\begin{equation}\label{eqn:R_N}
R_N:Y\to Y_N,\qquad R_N(A)\coloneqq P_N^2AP_N^1\,.
\end{equation}
Equip $Y_N$ with the Hilbert--Schmidt norm $\|\cdot\|_{\mathrm{HS}}$, which is independent of the choice of basis of $V_N^1$ and $V_N^2$.
Fix orthonormal bases $\{\varphi_j\}_{j=1}^{n_1}$ of $V_N^1$ and $\{\psi_i\}_{i=1}^{n_2}$ of $V_N^2$ and define the vectorization
\[
\mathrm{vec}:Y_N\to\R^{d_N},\qquad
\mathrm{vec}(A)\coloneqq \Big(\langle A\varphi_j,\psi_i\rangle_{V_N^2}\Big)_{\substack{1\le i\le n_2\\1\le j\le n_1}}\,,
\]
thus $\|\mathrm{vec}(A)\|_2=\|A\|_{\mathrm{HS}}$ for all $A\in Y_N$. By introducing the finite dimensional restriction, we have following standard inequality, which records that (i) restriction and projection cannot increase operator norm, and (ii) the Hilbert–Schmidt norm of the restricted operator is controlled by the original operator norm up to a factor $\sqrt{N}$. 
\begin{lemma}\label{lem:RN-stable}
For any $A\in Y$,
\[
\|R_N(A)\|_{\mathrm{op}(V_N^1\to V_N^2)}\ \le\ \|A\|_Y,
\qquad
\|R_N(A)\|_{\mathrm{HS}}\ \le\ \sqrt{N}\,\|A\|_Y\,.
\]
\end{lemma}

\begin{proof}[Proof of Lemma ~\ref{lem:RN-stable}]
Since $P_N^1,P_N^2$ are orthogonal projections, $\|P_N^1\|=\|P_N^2\|=1$, hence
\[
\|R_N(A)\|_{\mathrm{op}}=\|P_N^2AP_N^1\|_{\mathrm{op}}\le \|A\|_Y.
\]
For the Hilbert--Schmidt norm, for any linear map $B:V_N^1\to V_N^2$,
$\|B\|_{\mathrm{HS}}\le \sqrt{\mathrm{rank}(B)}\,\|B\|_{\mathrm{op}}$ and
$\mathrm{rank}(R_N(A))\le N$, so
\[
\|R_N(A)\|_{\mathrm{HS}}
\le \sqrt{N}\,\|R_N(A)\|_{\mathrm{op}}
\le \sqrt{N}\,\|A\|_Y.
\]
\end{proof}
Finally we define the discretized DtN measurement scaled by $1/\sqrt{N}$,
\begin{equation}\label{eqn:F_N}
F_N:K\to\R^{d_N},\qquad F_N(\gamma)\coloneqq \frac{1}{\sqrt{N}}\mathrm{vec}\big(R_N(\Lambda_\gamma)\big)\,,
\end{equation}
and establish the Chartwise Lipschitz for $F_N$ as follows. The $1 / \sqrt{N}$ normalization rescales distances in $\R^{d_N}$, and Lipschitz constants rescale accordingly. In particular, this normalization does not change entropy exponents (rates) and only rescales accuracy parameter $\varepsilon$. We adopt the $1/\sqrt{N}$ normalization to absorb the discretization-induced factor $\sqrt{N}$, which can streamline the subsequent Lipschitz and entropy bounds. We also discuss the effect of such normalization in Appendix~\ref{sec:normalization}.

We now combine the chartwise regularity of DtN mapping and Lemma~\ref{lem:RN-stable} to conclude that, on each parameter patch $E_i$, the chartwise discretized forward map $G_{N,i}$ is Lipschitz in the Euclidean norm.
\begin{lemma}[Chartwise Lipschitz for measurements]\label{lem:GN-Lip}
For each $i$, define $G_{N,i}:E_i\to\R^{d_N}$ by $G_{N,i}(v)\coloneqq F_N(\varphi_i^{-1}(v))$.
Then $G_{N,i}$ is Lipschitz on $E_i$ with
\[
\|G_{N,i}(v)-G_{N,i}(v')\|_2 \le \,M_i\,\|v-v'\|_2,\qquad \forall v,v'\in E_i\,,
\]
where $M_i\coloneqq \sup_{v\in\conv(E_i)}\|D\widetilde\Lambda_i(v)\|_{\mathrm{op}}<\infty$ and independent of $N$.

\end{lemma}

\begin{proof}[Proof of Lemma ~\ref{lem:GN-Lip}]
Using $\|\mathrm{vec}(\cdot)\|_2=\|\cdot\|_{\mathrm{HS}}$ and Lemma~\ref{lem:RN-stable},
\[
\|G_{N,i}(v)-G_{N,i}(v')\|_2
=\frac{1}{\sqrt{N}}\big\|R_N\big(\widetilde\Lambda_i(v)-\widetilde\Lambda_i(v')\big)\big\|_{\mathrm{HS}}
\le\ \|\widetilde\Lambda_i(v)-\widetilde\Lambda_i(v')\|_Y\,.
\]
Then it suffices to show
\[
\|\widetilde\Lambda_i(v)-\widetilde\Lambda_i(v')\|_Y\le M_i\|v-v'\|_2\,.
\]
To this end, we apply Fundamental theorem of calculus along the segment connecting $v$ and $v'$, 
\[
\widetilde\Lambda_i(v')-\widetilde\Lambda_i(v)=\int_0^1 D\widetilde\Lambda_i\big(v+t(v'-v)\big)\,(v'-v)\,dt\,.
\]
\[
\|\widetilde\Lambda_i(v)-\widetilde\Lambda_i(v')\|_Y\le \sup_{t\in[0,1]}\|D\widetilde\Lambda_i\big(v+t(v'-v)\big) \|_{\mathrm{op}}\|v-v'\|_2 \leq \sup_{v\in \operatorname{conv}(E_i)}\big\|D\widetilde\Lambda_i(v)\big\|_{\mathrm{op}} \|v-v'\|_2 \,.
\]
Here we apply triangle inequality for integrals and the definition of operator norm in the first inequality, and in the last inequality we use the facts that $\operatorname{conv}(E_i)\subset \mathcal O_i$ and $D\widetilde\Lambda_i$ is continuous on $\mathcal O_i$ by Lemma~\ref{lem:DtN-C1-chart}.


\end{proof}
\subsection{Metric entropy and bounded support of $F_N(K)$}\label{sec:entropy-bdd-FN}
The diffusion convergence theorem we invoke later depends on the intrinsic complexity of the data
distribution, quantified via covering numbers of its support at a prescribed resolution
$\varepsilon_0 = T^{-c_{\varepsilon_0}}$, where $T\in\mathbb{N}$ is the number of diffusion steps and
$c_{\varepsilon_0}>0$ is fixed.
We therefore aim to bound the metric entropy of the measurement set $F_N(K)$ with respect to the
Euclidean norm $\|\cdot\|_2$.

\begin{definition}[Covering number and metric entropy]\label{def:covering}
Let $(\mathcal{X},d)$ be a metric space and $S\subset\mathcal{X}$.
For $\varepsilon>0$, the covering number $\mathcal{N}_\varepsilon(S;d)$ is the smallest $n$ such that
$S\subset\bigcup_{j=1}^n B(x_j,\varepsilon)$ for some $x_1,\dots,x_n\in\mathcal{X}$, where
$B(x,\varepsilon)\coloneqq \{y:\ d(y,x)\le\varepsilon\}$.
The metric entropy is defined as $\log \mathcal{N}_\varepsilon(S;d)$.
\end{definition}

The following lemma shows that an \(L\)-Lipschitz map does not increase covering numbers after rescaling
the covering radius.

\begin{lemma}[Coverings under Lipschitz maps]\label{lem:cover-lip}
Let $(\mathcal X,d_{\mathcal X})$ and $(\mathcal Z,d_{\mathcal{Z}})$ be metric spaces.
If $f:\mathcal X\to\mathcal Z$ is $L$-Lipschitz, then for any $S\subset\mathcal X$ and $\varepsilon>0$,
\[
\mathcal{N}_\varepsilon\!\bigl(f(S);d_{\mathcal{Z}}\bigr)\ \le\ \mathcal{N}_{\varepsilon/L}\!\bigl(S;d_{\mathcal X}\bigr)\,.
\]
\end{lemma}

\begin{proof}[Proof of Lemma ~\ref{lem:cover-lip}]
Let $\{x_j\}_{j=1}^n$ be an $\varepsilon/L$-net of $S$ in $d_{\mathcal X}$ with
$n=\mathcal{N}_{\varepsilon/L}\!\bigl(S;d_{\mathcal X}\bigr)$, so
$S\subset\bigcup_{j=1}^n B_{\mathcal X}(x_j,\varepsilon/L)$.
For any $z=f(s)\in f(S)$ choose $j$ with $d_{\mathcal X}(s,x_j)\le\varepsilon/L$.
Then $d_{\mathcal{Z}}(z,f(x_j))\le L d_{\mathcal X}(s,x_j)\le\varepsilon$.
This implies that $\{f(x_j)\}$ is an $\varepsilon$-net of $f(S)$, and as a consequence,
the covering number of $f(S)$ is no greater than $n$.
\end{proof}

\begin{remark}
If a map is rescaled by a scalar factor $a>0$ (e.g.\ $a=1/\sqrt{N}$ in the definition of $F_N$),
then $a f$ is $(aL)$-Lipschitz whenever $f$ is $L$-Lipschitz. Thus Lemma~\ref{lem:cover-lip}
shows that such normalizations only rescale the covering radius (hence modify constants), and do not
affect the exponent in the metric-entropy bounds.
\end{remark}

We also need a volumetric covering bound in $\R^{2n_{\mathrm v}}$, since an $n_{\mathrm v}$-gon in $\R^2$
is parameterized by its $n_{\mathrm v}$ ordered vertices in $\R^2$, i.e.\ $v\in\R^{2n_{\mathrm v}}$.

\begin{lemma}[Volumetric covering in $\R^{2n_{\mathrm v}}$]\label{lem:volumetric}
There exists an absolute constant $C_{\mathrm{vol}}>0$ such that the following holds.
If $E\subset\R^{2n_{\mathrm v}}$ satisfies $E\subset B_2(0,R)$ for some $R>0$, then for every $\delta>0$,
\[
\mathcal{N}_\delta(E;\|\cdot\|_2)\ \le\ \max\Big\{\,1,\ \Big(\frac{C_{\mathrm{vol}}\,R}{\delta}\Big)^{2n_{\mathrm v}}\Big\}\,.
\]
\end{lemma}


\begin{remark}
Applying Lemma~\ref{lem:volumetric} to a globally rescaled set $aE$ simply replaces $R$ by $aR$, where $a>0$ is an arbitrary constant. Consequently, the $1/\sqrt{N}$ normalization in $F_N$ can only enter through prefactors such as
$C_{\mathrm{vol}}R$ (or Lipschitz constants), while the exponent $2n_{\mathrm v}$ remains unchanged.
\end{remark}

We now bound (i) the metric entropy of $F_N(K)$ and (ii) a uniform radius bound for its support.
The proof covers $K$ by finitely many patches, uses Lipschitz bounds for the chartwise measurement maps,
and applies a volumetric covering estimate in $\R^{2n_{\mathrm v}}$.

\begin{lemma}[Uniform boundedness of DtN operator on $K$]\label{lem:DtN-bdd}
Define \(r_K \coloneqq  \sup_{\gamma\in K}\|\Lambda_\gamma\|_Y\). Then \(r_K<\infty\).
\end{lemma}

\begin{proof}[Proof of Lemma ~\ref{lem:DtN-bdd}]
On $K$, $\gamma$ is uniformly elliptic and bounded by
$\min\{1,\kappa\}\le \gamma(x)\le \max\{1,\kappa\}$.
Standard energy estimates for elliptic boundary value problems imply that
$\|\Lambda_\gamma\|_Y$ is uniformly bounded over $\gamma\in K$, hence $r_K<\infty$ \citep{evans2022partial}.
\end{proof}

We could now define the maximum radius for all local coordinate patches
\begin{equation}\label{eq:Rstar_poly}
R_*\coloneqq \max_{1\le i\le m}\ \sup_{v\in E_i}\|v\|_2<\infty.
\end{equation}
\begin{theorem}[Metric entropy and bounded support of discretized DtN measurements]\label{thm:FN-main}
Consider $F_N:K\to\R^{d_N}$ in \Eqref{eqn:F_N}. Fix $c_{\varepsilon_0}>0$ and for all $T\ge 2$ set
$\varepsilon_0\coloneqq T^{-c_{\varepsilon_0}}$.
Let $C_{\mathrm{vol}}>0$ be as in Lemma~\ref{lem:volumetric}. Then for all $T\ge 2$,
\begin{equation}\label{eq:entropy-FNK}
\log \mathcal{N}_{\varepsilon_0}\big(F_N(K);\|\cdot\|_2\big)\ \le\ C_F\,\log T \,,
\end{equation}
where one may take
\begin{equation}\label{eq:CF}
C_F \ \coloneqq \ 2n_{\mathrm v}\,c_{\varepsilon_0}\ +\ \frac{\log m + 2n_{\mathrm v}\log\!\big(C_{\mathrm{vol}}\,R_*\,M_*\big)}{\log 2}\,.
\end{equation}
Here $m$ and $R_*$ are as in Lemma~\ref{lem:finite-atlas}, \(M_i\) is the Lipschitz constant of \(G_{N,i}\)
given in Lemma~\ref{lem:GN-Lip}, and \(M_*\coloneqq \max_{1\le i\le m} M_i\).
Moreover, $F_N(K)$ is bounded in $\R^{d_N}$ and satisfies
\[
\sup_{\gamma\in K}\|F_N(\gamma)\|_2 \le r_{K}\,,
\]
where \(r_K\) is the finite constant from Lemma~\ref{lem:DtN-bdd}.
\end{theorem}

\begin{proof}
By Lemma~\ref{lem:finite-atlas}, $K\subset\bigcup_{i=1}^m V_{P_i}$, hence
\[
F_N(K)\subset \bigcup_{i=1}^m F_N\big(K\cap \overline{V}_{P_i}\big)
= \bigcup_{i=1}^m G_{N,i}(E_i),
\]
where $E_i\subset\R^{2n_{\mathrm v}}$ is the $i$-th parameter patch and $G_{N,i}:E_i\to\R^{d_N}$ is defined
in Lemma~\ref{lem:GN-Lip}.
Therefore for any $\varepsilon>0$,
\begin{equation}\label{eqn:NF_N(K)}
\mathcal{N}_\varepsilon\big(F_N(K);\|\cdot\|_2\big)
\le \sum_{i=1}^m \mathcal{N}_\varepsilon\big(G_{N,i}(E_i);\|\cdot\|_2\big)
\le \sum_{i=1}^m \mathcal{N}_{\varepsilon/M_i}(E_i;\|\cdot\|_2),
\end{equation}
where the last inequality follows from Lemma~\ref{lem:GN-Lip} and Lemma~\ref{lem:cover-lip}.

Each $E_i\subset\R^{2n_{\mathrm v}}$ is bounded and satisfies $E_i\subset B_2(0,R_*)$ by definition of $R_*$.
Applying Lemma~\ref{lem:volumetric} gives, for every $\varepsilon>0$,
\[
\mathcal{N}_{\varepsilon/M_i}(E_i;\|\cdot\|_2)
\le \max\Big\{1,\ \Big(\frac{C_{\mathrm{vol}}\,R_*\,M_i}{\varepsilon}\Big)^{2n_{\mathrm v}}\Big\}.
\]
Hence,
\[
\mathcal{N}_\varepsilon\big(F_N(K);\|\cdot\|_2\big)
\le m\max\Big\{1,\ \Big(\frac{C_{\mathrm{vol}}\,R_*\,M_*}{\varepsilon}\Big)^{2n_{\mathrm v}}\Big\}.
\]
Taking $\varepsilon=\varepsilon_0=T^{-c_{\varepsilon_0}}$ yields
\[
\log \mathcal{N}_{\varepsilon_0}\big(F_N(K)\big)
\le \log m + 2n_{\mathrm v}\log\!\big(C_{\mathrm{vol}}R_*M_*\big) + 2n_{\mathrm v}c_{\varepsilon_0}\log T.
\]
Since $T\ge 2$, the additive constant satisfies
\[
\log m + 2n_{\mathrm v}\log\!\big(C_{\mathrm{vol}}R_*M_*\big)
\le \frac{\log m + 2n_{\mathrm v}\log\!\big(C_{\mathrm{vol}}R_*M_*\big)}{\log 2}\,\log T,
\]
which gives \Eqref{eq:entropy-FNK} with $C_F$ in \Eqref{eq:CF}.

We now show a uniform radius bound for $F_N(K)$.
Using $\|F_N(\gamma)\|_2=\frac{1}{\sqrt{N}}\|R_N(\Lambda_\gamma)\|_{\mathrm{HS}}$,
Lemma~\ref{lem:RN-stable}, and Lemma~\ref{lem:DtN-bdd},
\[
\|F_N(\gamma)\|_2 \le \|\Lambda_\gamma\|_Y \le r_{K}.
\]
Taking the supremum over $\gamma\in K$ completes the proof.
\end{proof}

\subsection{Convergence guarantee for diffusion completion}\label{sec:ddpm_FNK}
Consider the (scaled) full discretized DtN measurement vector
\[
X_0 \coloneqq  F_N(\gamma)\in\R^{d_N},
\]
and the corresponding observed data
\[
Z \coloneqq  (\mathbf{M},\; \mathbf{M}\odot X_0)\in \{0,1\}^{d_N}\times\R^{d_N}.
\]
Note that $F_N$ includes the normalization $1/\sqrt{N}$ in \Eqref{eqn:F_N}.

We recall the unconditional result of \citet[Thm.~2]{liang2025low} and will apply it to each conditional
distribution $p_{X_0\mid Z=z}$ by treating $z$ as a fixed parameter. 

\begin{lemma}[Unconditional DDPM convergence under low intrinsic complexity
{\rm(\citealp[Thm.~2]{liang2025low})}]\label{lem:liang-ddpm}
Let $X_0\sim p_{\mathrm{data}}$ be a distribution on $\R^d$ satisfying Assumptions~1--3 of
\citet{liang2025low} (intrinsic dimension $k$ quantified by metric entropy at $\epsilon_0=T^{-c_{\epsilon_0}}$,
bounded support with radius $R=T^{c_R}$, and averaged $\ell_2$ score error $\varepsilon_{\mathrm{score}}$).
Then the DDPM sampler satisfies
\begin{equation}\label{eq:TV_bound_uncond}
\mathrm{TV}\bigl(p_{X_1},\,p_{Y_1}\bigr)
\ \lesssim\ \frac{k\,\log^3 T}{T}\ +\ \varepsilon_{\mathrm{score}}\sqrt{\log T}\,.
\end{equation}
\end{lemma}

In our polygonal EIT setting, the natural parameter dimension is $k=2n_{\mathrm v}$ (ordered vertices in
$\R^2$). The metric entropy bound for intrinsic dimension and the boundedness of support have been verified in
Theorem~\ref{thm:FN-main}, which depends on the assumptions for admissible polygons. We assume the conditional score estimation error as follows.

\begin{assumption}[$\ell_2$ conditional score error]\label{ass:cond_score}
There exists $\varepsilon_{\mathrm{score}}\ge 0$ such that for $Z$-almost every $z$,
\[
\frac{1}{T}\sum_{t=1}^T
\mathbb E\!\left[\ \big\|s_{\theta}(X_t,z,t)-s_t^{\star}(X_t,z)\big\|_2^2\ \mid Z=z \right]
\ \le\ \varepsilon_{\mathrm{score}}^2\,.
\]
\end{assumption}

\begin{remark}\label{rem:score-assump}
Assumption~\ref{ass:cond_score} abstracts the error incurred by learning the conditional score from data
using a neural network. Providing a rigorous, nonasymptotic bound for $\varepsilon_{\mathrm{score}}$ is a
nontrivial learning-theoretic problem and typically depends on the model class, training procedure, and
regularity properties of the target conditional score.
In this paper we follow the common approach in the diffusion convergence literature and treat
$\varepsilon_{\mathrm{score}}$ as an exogenous accuracy parameter, as in
\citep{chen2023improved,benton2023nearly,liang2025low}.
\end{remark}
With conditional score error estimation in Assumption~\ref{ass:cond_score}, we
verify the assumptions of \citet[Thm.~2]{liang2025low} for each conditional law $p_{X_0\mid Z=z}$ and obtain
the following guarantee.
\begin{theorem}[Conditional DDPM completion error in TV]\label{thm:cond-ddpm}
Consider the set of full DtN measurements from all admissible conductivities $F_N(K)\subset\R^{d_N}$, and let
$\epsilon_0\coloneqq T^{-c_{\epsilon_0}}$ with $c_{\epsilon_0}>0$ as in \citealp[Assumption~1]{liang2025low}. Assume furthermore that Assumption~\ref{ass:cond_score} holds. Then for $Z$-almost every $z$, the conditional DDPM sampler produces $Y_1 \mid Z=z$ whose law satisfies
\begin{equation}\label{eq:TV_bound_cond}
\mathrm{TV}\bigl(p_{X_1\mid Z=z},\,p_{Y_1\mid Z=z}\bigr)
\ \lesssim\ \frac{2n_{\mathrm v}\,\log^3 T}{T}\ +\ \varepsilon_{\mathrm{score}}\sqrt{\log T}\,.
\end{equation}
\end{theorem}

\begin{proof}
Fix $z$ such that Assumption~\ref{ass:cond_score} holds. Since $X_0\in F_N(K)$ almost surely, it follows
that $\mathrm{supp}(p_{X_0\mid Z=z})\subset F_N(K)$. Consequently, for any $\epsilon>0$,
\[
\mathcal{N}_\epsilon\bigl(\mathrm{supp}(p_{X_0\mid Z=z});\|\cdot\|_2\bigr)
\ \le\ \mathcal{N}_\epsilon\bigl(F_N(K);\|\cdot\|_2\bigr),
\qquad
\sup_{x\in\mathrm{supp}(p_{X_0\mid Z=z})}\|x\|_2\ \le\ r_K.
\]

Taking $\epsilon=\epsilon_0$ yields
\[
\log \mathcal{N}_{\epsilon_0}\bigl(\mathrm{supp}(p_{X_0\mid Z=z})\bigr)\ \le\ C_F\log T,
\]
by metric entropy bound from Theorem~\ref{thm:FN-main}.
Therefore \citealp[Assumption~1]{liang2025low} holds for $p_{X_0\mid Z=z}$ with intrinsic dimension
$k=2n_{\mathrm v}$ and constant $C_{\mathrm{cover}}\coloneqq C_F/(2n_{\mathrm v})$.
The radius bound from Theorem~\ref{thm:FN-main} verifies \citealp[Assumption~2]{liang2025low}, since $r_K$ does not depend on $T$ and we can choose $c_R\ge \log_2 r_K$ such that $r_K\le T^{c_R}$ for $T\ge 2$.

Assumption~\ref{ass:cond_score} is the conditional analogue of \citealp[Assumption~3]{liang2025low}.
For each fixed $z$, the conditional DDPM is an unconditional DDPM on $\R^{d_N}$ with data law
$p_{X_0\mid Z=z}$. Applying Lemma~\ref{lem:liang-ddpm} to $p_{X_0\mid Z=z}$ yields \Eqref{eq:TV_bound_cond}.
\end{proof}

\begin{remark}[Why the conclusion is stated at time $t=1$]
The TV guarantee above compares $p_{X_1\mid Z=z}$ and $p_{Y_1\mid Z=z}$, consistent with
\citealp[Thm.~2]{liang2025low}. For data distributions supported on low-dimensional sets in $\R^d$,
a TV guarantee directly for $p_{X_0\mid Z=z}$ is generally impossible against any absolutely continuous
sampler output, due to mutual singularity.
\end{remark}

The ambient dimension of the measurement vector is $d_N=N^2$, whereas the admissible family $F_N(K)$ is
parameterized by only $2n_{\mathrm v}$ real parameters (ordered vertices in $\R^2$).
This low-dimensional structure is reflected in the metric-entropy bound \Eqref{eq:entropy-FNK}:
after passing to local coordinates, the forward map is Lipschitz on each patch, so the covering numbers of
$F_N(K)$ are controlled by volumetric coverings in $\R^{2n_{\mathrm v}}$, yielding the power $2n_{\mathrm v}$
in the covering-number bound.

\section{Numerical Experiments}\label{sec:numeric}

In this section, we evaluate the proposed method against baseline approaches on two tasks: (i) completing full DtN measurements from randomly sampled sparse measurements, and (ii) reconstructing EIT conductivities from sparse measurements. The details of the experimental setup and model training are provided in Sections~\ref{sec:data_gen} and~\ref{sec:Implement_details}, while the testing results are reported in Sections \ref{sec:result_DC} and~\ref{sec:result_IP}.

\subsection{Data generation}\label{sec:data_gen}
\paragraph{Discretizations} We consider the unit disk as the computational domain $\Omega$, which is discretized over a conforming triangular mesh (see Figure \ref{fig:conductivity_discretizations}). 
 This results in $N_\mathcal{I} = 1324$ interior nodal points, $N_\mathcal{B} = 128$ boundary nodal points, and 2774 triangular elements. The elliptic PDE \Eqref{eq:forward} is then solved using the finite element method (FEM) with P1 elements. Consequently, the DtN measurement matrix has size $128 \times 128$, while the conductivity $\gamma$ is discretized into a vector of length 2774 at the centroids of all elements. To facilitate the use of conductivity data in deep learning models, $\gamma$ is zero-padded and extended to a function on the square $[-1,1]\times[-1,1]$, and then discretized over a $128 \times 128$ grid via linear interpolation, as shown in Figure~\ref{fig:conductivity_discretizations}.  

\paragraph{Conductivity distributions}
We consider two different distributions of conductivity throughout this section, which are defined as follows.
\begin{itemize}
\item \textbf{Disks:} 
This dataset consists of conductivities over a unit disk $\Omega$ with unit background value, together with constant values over a few circular inclusions. In particular, the conductivities are defined as the following.
\[
\gamma(x) \;=\; 1 \;+\; \sum_{i=1}^{n} \bigl(c_i-1\bigr)\,\mathbf{1}_{\{\|x-x_i\|_2 \le r_i\}}(x), \qquad x\in\Omega\,,
\]
where \(n\in\{2,3,4,5\}\), and the inclusion centers \(x_i\in\Omega\), radii \(r_i\), and contrasts \(c_i\) vary across samples.
For each fixed \(n\), we independently draw radii $r_i \sim \mathcal{U}[0.2,\,0.4]$, contrasts $c_i \sim \mathcal{U}[2,\,8]$, and centers $x_i$ uniformly over $\Omega$. We use rejection sampling to ensure that the inclusions are fully contained within $\Omega$ and mutually disjoint. That is,
\[
\dist(x_i,\partial\Omega)\;>\; r_i \quad\text{for all } i,
\qquad\text{and}\qquad
\|x_i-x_j\|_2 \;>\; r_i+r_j \quad\text{for all } i\neq j\,.
\]
For each \(n\in\{2,3,4,5\}\), \(500\) admissible training samples and \(100\) admissible test samples are generated, yielding a total of \(2000\) training and \(400\) test conductivities.

\item \textbf{Shepp--Logan Phantom:} This dataset is based on \citet{shepp1974fourier}, a standard benchmark in medical imaging that consists of piecewise constant functions that model human skulls, ventricles and tumors.
Our code follows the setup in \citet{Chung2020}.
\end{itemize}
Examples of conductivities in each dataset can be found in the leftmost columns of Figure~\ref{fig:disk_deepsolver_noiseless} and Figure~\ref{fig:sl_deepsolver_noiseless} respectively.

\begin{figure}[!htbp]
    \centering
    \includegraphics[width=0.8\linewidth]{discretization.png}
\caption{Discretization of conductivity $\gamma$. \textbf{Left:} Triangular mesh over the unit disk where red boundary points mark candidate sensor locations. 
\textbf{Middle:} A conductivity over the triangular mesh from the Disk distribution. 
\textbf{Right:} The same conductivity converted to a uniform $128\times128$ Cartesian grid on $[-1,1]\times[-1,1]$ with zero padding.
}
\label{fig:conductivity_discretizations}
\end{figure}

\paragraph{Normalization of the DtN measurement} 
The DtN measurement $\mathbf{\Lambda}_\gamma$ exhibits strong heterogeneity between diagonal and off-diagonal entries, as the off-diagonal entries are fast-decaying for any $L^\infty$ coefficient $\gamma$~\citep{bebendorf2003existence,bechtel2024off}. 
Furthermore, the ill-posedness of the Calder\'on problem suggests that the difference between two measurements $\mathbf{\Lambda}_{\gamma_1}$ and $\mathbf{\Lambda}_{\gamma_2}$ is highly insensitive to their corresponding conductivity differences. As a consequence, all DtN measurements are similar to those of the background (unit) conductivity~\citep{alessandrini1988stable}. Such similarity in the data causes a significant challenge in training effective deep learning models. 
To this end, we normalize each DtN measurement by that of the background (unit) conductivity.
The normalized DtN measurement is defined as the elementwise (Hadamard) ratio $\bar{\mathbf{\Lambda}}_\gamma\coloneqq \mathbf{\Lambda}_\gamma \oslash \mathbf{\Lambda}_{\textbf{1}}$. 
All models are trained and evaluated on the normalized DtN measurement \(\bar{\mathbf{\Lambda}}_\gamma\). As illustrated in Figure~\ref{fig:DtN_normalization}, raw DtN measurements are visually similar across distinct conductivities, whereas their normalized version reveals discriminative patterns.


\begin{figure}[!htbp]
    \centering
    \includegraphics[width=0.7\linewidth]{DtN_normalization.png}
\caption{Normalization of DtN measurements. 
\textbf{Left column:} conductivity fields (top: background; middle: disks; bottom: Shepp--Logan phantom). 
\textbf{Middle column:} raw DtN matrices computed from the corresponding conductivities. 
\textbf{Right column:} normalized DtN matrices produced by our preprocessing step.}
\label{fig:DtN_normalization}
\end{figure}



\paragraph{Measurement configurations} 
Two types of mask matrices are considered in our numerical experiments.
\begin{itemize}
    \item \textbf{Principal submatrix}: 
    Sources and receivers are placed on a fixed number of boundary nodes that are chosen uniformly at random. In this case, the mask matrix $\mathbf{M}_s$ is a principal submatrix, i.e.
\[
(\mathbf{M}_{s})_{ij} = 
\begin{cases}
1, & i\in\mathcal{S} \text{ and } j\in\mathcal{S}\,,\\[2pt]
0, & \text{otherwise}\,.
\end{cases}
\]
    where an index set $\mathcal{S} \subset [N_{\mathcal B}]$ is drawn uniformly at random without replacement. The total number of boundary nodes $|\mathcal{S}| = \lfloor \sqrt{s}\,N_{\mathcal B}\rfloor$ is determined by the sampling rate $s\in (0,1)$. 
\item \textbf{Random}: Each entry of $\mathbf{\Lambda}_\gamma$ is observed independently with probability $s$ as in \citet{bui2022bridging}. The mask $\mathbf{M}_s$ is therefore a Bernoulli random matrix with parameter $s$, i.e.,
\[
(\mathbf{M}_{s})_{ij} = 
\begin{cases}
1, & \text{w.p. } s\,,\\[2pt]
0, & \text{w.p. } 1-s\,.
\end{cases}
\]
\end{itemize}
 
Examples of both types of mask matrices are plotted in Figure \ref{fig:comparison_completion_DM_MC}.

\paragraph{Noise injection} To investigate the robustness of our methods, we apply a noise pattern considered in \citet{borcea2013study}, which is similar to multiplicative noise, to both training and testing measurements.
\begin{equation*}
    \mathbf{\Lambda}_{\gamma}^{\text{Noisy}} = \mathbf{\Lambda}_{\gamma} + \sigma_\text{noise}\mathbf{\Lambda}_{\textbf{1}}\odot\mathcal{E} \,,
\end{equation*}
where $\mathbf{\Lambda}_{\textbf{1}}$ is the background DtN measurement, $\sigma_\text{noise}\in (0,1)$ is the strength of noise, and $\mathcal{E} \in \mathbb{R}^{N_{\mathcal{B}}\times N_{\mathcal{B}}}$ is a random matrix whose entries are independently drawn from the standard normal distribution. We consider $\sigma_\text{noise}= 0.05 (5\%)$ in our experiments, and inject noise to $\mathbf{\Lambda}_\gamma$ first and then apply the normalization to obtain $\bar{\mathbf{\Lambda}}_\gamma$.


\subsection{Implementation details}\label{sec:Implement_details}
We evaluate the proposed method through two stages: data completion in Section \ref{sec:result_DC} and conductivity reconstruction in Section \ref{sec:result_IP}. 
In the data completion stage, we compare our method with the hierarchical completion baseline introduced in~\citet{bui2022bridging}. In the conductivity reconstruction stage, we adopt \texttt{IAE-Net}, an efficient neural operator introduced in~\citet{ong2022integral}, which achieves excellent accuracy on forward and inverse PDE operator learning tasks.
To demonstrate the effectiveness of our proposed method, we compare the results of three inverse problem solving strategies illustrated in Figure~\ref{fig:diagram}: 
(i) \emph{direct reconstruction from sparse measurements}; 
(ii) \emph{data completion followed by reconstruction}, where the data completion part is implemented by our diffusion model or by the hierarchical matrix completion; and (iii) \emph{direct reconstruction from full measurements}, which sets the best achievable accuracy for the specified inverse solver (e.g. \texttt{IAE-Net}).
We next details the implementations of our diffusion completion model in Section~\ref{sec:Implement_details}, the hierarchical matrix completion baseline, and the \texttt{IAE-Net} inverse solver, along with the evaluation metrics. 

\subsubsection{Data completion methods}
\paragraph{Diffusion model implementation}
We adopt the Elucidated Diffusion Model (EDM) framework \citep{karras2022elucidating} with the preconditioned denoising network $D_{\theta}$, which can be interpreted as a reparameterization of the score functions studied in \citet{song2020score, ho2020denoising}. The denoiser network $D_\theta(x,\sigma)$ in \citet{karras2022elucidating} predicts a clean DtN measurement $x_0$ from a perturbed input $x = x_0 + \sigma \varepsilon$, where $\varepsilon \sim \mathcal{N}(0,I)$ and $\sigma$ is the noise level. 
For training, we use the same loss function and noise-level sampling strategy as in \citet{karras2022elucidating}, and train the model on normalized DtN measurement data. 
The denoiser $D_{\theta}$ is implemented as a 2D U-Net \citep{ronneberger2015u}, where each DtN matrix is treated as a single-channel image. The network has a base width of 128, includes two residual blocks at each resolution level, and performs dyadic downsampling and upsampling by a factor of 2, repeated three times, with an encoder--decoder symmetric structure and skip connections. Each residual block consists of $3\times 3$ convolutions, normalization, and SiLU activations. The scalar noise level $\sigma$ is encoded using a Fourier embedding and injected into the residual blocks via affine modulation. Attention mechanisms are applied only in the bottleneck layer. 
Optimization is carried out with AdamW \citep{loshchilov2017decoupled} using a learning rate of $10^{-4}$, an exponential moving average of model parameters with decay $0.999$, and a batch size of $32$. Unless otherwise specified, all remaining EDM hyperparameters and conventions follow \citet{karras2022elucidating}.

\paragraph{Baseline model: hierarchical matrix completion} 
\begin{figure}[!htbp]
    \centering
    \includegraphics[width=1\linewidth]{h_matrix_visualization.png}
    \caption{Hierarchical partition with level 3. \emph{Diagonal blocks} are colored in orange, and \emph{off-diagonal blocks} are colored in blue.}
    \label{fig:h-matrix}
\end{figure}
We consider the hierarchical matrix completion method in~\citet{bui2022bridging} as a baseline for completing DtN measurements. This method exploits the hierarchical off-diagonally low-rank structure of the DtN measurement matrix \citep{bebendorf2003existence,ballani2017matrices} and employs standard matrix completion methods~\citep{recht2011simpler}. Concretely, the DtN matrix is divided into a hierarchical tree of diagonal blocks and off-diagonal blocks as shown in Figure \ref{fig:h-matrix}. 
To apply matrix completion globally, a corresponding hierarchical masking matrix needs to be considered for this method.
In particular, it is assumed that the observed EIT measurements consist of all data on the diagonal blocks and partial data that are sampled uniformly at random over all off-diagonal blocks. See Figure \ref{fig:mask_comparison} for the global masking matrix considered in such hierarchical setting. 
Specifically, we implement a three-level decomposition of the measurement matrix, which leads to in total eight diagonal blocks and fourteen off-diagonal blocks. 
The matrix completion algorithm is applied in parallel to each of the off-diagonal blocks to obtain the completed data, which are then assembled together to form the completed measurement matrix.
In particular, consider an arbitrary off-diagonal block $\mathbf{\Lambda}_{\gamma}^{l} \in \mathbb{R}^{\frac{N_{\mathcal B}}{2^l}\times \frac{N_{\mathcal B}}{2^l}}$ at level $l=1,2,3$ and a random mask $\mathbf{M}_s^{l}$ of the same size, the matrix completion algorithm outputs a solution to the following constrained nuclear norm optimization problem
\begin{equation}
\label{eq:mc_nucnorm}
\widehat{\mathbf{\Lambda}}_\gamma^{l}
\;=\;
\arg\min_{X} \ \|X\|_{*}
\quad \text{subject to} \quad
\mathbf{M}_s^{l}\odot X \;=\; \mathbf{M}_s^{l}\odot \mathbf{\Lambda}_\gamma^{l}\,, \quad X \in\mathbb{R}^{\frac{N_{\mathcal B}}{2^l}\times \frac{N_{\mathcal B}}{2^l}}\,.
\end{equation}
where \(\|\cdot\|_{*}\) denotes the matrix nuclear norm. 
The above minimization problem admits a unique solution if the sampling rate in the random mask $\mathbf{M}_s^{l}$ is sufficiently large.
We implement the nuclear optimization \Eqref{eq:mc_nucnorm} with the \texttt{MOSEK} solver from the \texttt{CVXPY} package.

\subsubsection{Inverse problem solver}  
We employ the neural operator \texttt{IAE-Net}~\citep{ong2022integral} as our inverse solver, for its excellent accuracy across various PDE forward and inverse operator learning tasks.
Two $\texttt{IAE-Net}$ models are trained separately on clean data to serve as a full-data solver and sparse-data solver respectively. 
In particular, the full-data solver is trained to learn the operator \(\mathcal{G}:\bar{\mathbf{\Lambda}}_\gamma \mapsto \gamma\) that maps the full DtN measurements to target conductivity field, while the sparse-data solver approximates \(\mathcal{G}_{\mathrm{masked}}: \bar{\mathbf{\Lambda}}_\gamma^{\text{o}} \mapsto \gamma\), where \(\bar{\mathbf{\Lambda}}_\gamma^{\text{o}} = \mathbf{M}_s \odot \bar{\mathbf{\Lambda}}_\gamma\). Specifically, $\mathbf{M}_s$ is a principal submatrix mask resampled independently at each training iteration, enabling the model to adapt to diverse sparse-measurement patterns. 

Both \texttt{IAE-Net} models are configured with 32 modes and a stack of IAE blocks, using four blocks for the Disks dataset and two for the Shepp--Logan dataset. Training minimizes the mean-squared error between predicted and ground-truth conductivities, optimized with Adam and weight decay $10^{-4}$. The learning rate follows a reduce-on-plateau schedule (factor $0.5$, patience $20$), with early stopping after $40$ stagnant epochs and a maximum of $300$ epochs. The initial learning rate is $10^{-2}$ for the Disks dataset and $10^{-3}$ for the Shepp--Logan dataset.

\subsubsection{Evaluation metrics}\label{sec:eval_metr}

In order to compare the proposed data completion method with the baseline, we consider the following evaluation metrics. These metrics quantify the fidelity of the completed DtN data and the accuracy of reconstructed conductivity respectively.

\paragraph{DtN data completion error}
A core objective of our framework is to accurately reconstruct the missing DtN measurements in $\mathbf{\Lambda}^{\text{o}}_\gamma$. We measure reconstruction quality using the normalized relative Frobenius error:
\begin{equation}
\label{eqn:dtn-error}
\text{RE}\left(\widehat{\bar{\mathbf{\Lambda}}}_\gamma\right) = \frac{\|\widehat{\bar{\mathbf{\Lambda}}}_\gamma  - \bar{\mathbf{\Lambda}}_\gamma \|_F}{\|\bar{\mathbf{\Lambda}}_\gamma \|_F}\,,
\end{equation}
where $\bar{\mathbf{\Lambda}}_\gamma$ denotes the normalized ground-truth DtN measurements and $\widehat{\bar{\mathbf{\Lambda}}}_\gamma$ is the completed normalized DtN measurements. 


\paragraph{Conductivity reconstruction error}
Note that the conductivity is defined on unit disk $\Omega \subset [-1,1]^2$ and padded zeros outside $\Omega$. We only measure the errors within $\Omega$. In particular, let \(\{x_i\}_{i=1}^{128}\) and \(\{y_j\}_{j=1}^{128}\) be the uniform grid coordinates on \([-1,1]\) and define
\[
\Omega_{\mathrm{grid}} \;=\; \bigl\{(i,j)\,:\ x_i^2 + y_j^2 \le 1 \bigr\},
\qquad
N_\Omega \;=\; \lvert \Omega_{\mathrm{grid}}\rvert\,.
\]
For a given neural network prediction $\hat{\gamma}$, we measure its  relative \(\ell_2\) error and the mean absolute error with respect to the ground truth $\gamma$ restricted to \(\Omega\):
\[
\mathrm{RE}\!\left(\hat{\gamma}\right)
= \frac{\displaystyle\sqrt{\sum_{(i,j)\in \Omega_{\mathrm{grid}}}\bigl(\hat{\gamma}_{i,j}-\gamma_{i,j}\bigr)^2}}
       {\displaystyle \sqrt{\sum_{(i,j)\in \Omega_{\mathrm{grid}}} \gamma_{i,j}^2}}, 
\qquad
\mathrm{MAE}\!\left(\hat{\gamma}\right)
= \frac{1}{N_\Omega}\sum_{(i,j)\in \Omega_{\mathrm{grid}}}\bigl|\hat{\gamma}_{i,j}-\gamma_{i,j}\bigr|\,.
\]

In addition, we also introduce structural similarity (SSIM) for evaluating conductivity reconstruction.
SSIM quantifies perceptual similarity by jointly comparing local luminance, contrast, and structure, thereby capturing preservation of geometry and edges more faithfully than pixel-wise metrics such as $\ell_2$ error or MAE.
Given two conductivities \(\gamma_{1},\gamma_2 \), the local structural similarity at pixel \(p\) is
\[
\mathrm{SSIM}_w(\gamma_1,\gamma_2)=
\frac{\bigl(2\mu_1(p)\mu_2(p)+C_1\bigr)\bigl(2\sigma_{12}(p)+C_2\bigr)}
     {\bigl(\mu_1(p)^2+\mu_2(p)^2+C_1\bigr)\bigl(\sigma_1^2(p)+\sigma_2^2(p)+C_2\bigr)}\,,
\]
where \(\mu_1,\mu_2\) are local means, \(\sigma_1^2,\sigma_2^2\) local variances, and \(\sigma_{12}\) the local covariance computed with a sliding window $w$. 
The reported image-level SSIM is the mean of these local SSIM:
\[
\mathrm{SSIM}(\gamma_1,\gamma_2)=\frac{1}{|\mathcal{W}|}\sum_{w\in\mathcal{W}}\mathrm{SSIM}_w(\gamma_1,\gamma_2)\,,
\]
with $\mathcal{W}$ the set of all sliding windows. In our implementation, we compute SSIM using the \texttt{scikit-image} package \citep{van2014scikit} with default settings.

\subsection{Results for data completion}\label{sec:result_DC}
We compare our diffusion-based data completion to the hierarchical matrix completion method as in \citet{bui2022bridging}. Since identical matrix completion algorithms \Eqref{eq:mc_nucnorm} are applied to all off-diagonal blocks with the same sampling rate, it is equivalent to compare data completion results restricted to a single off-diagonal block.
To this end, we implement our method and matrix completion over the upper right off-diagonal block at the first level in the subsequent subsection, and present the accuracy of data completion with various sampling rates from $s=1\%$ to $s=30\%$.

\paragraph{Off-diagonal block completion}
We evaluate both methods on the upper-right off-diagonal block of the DtN matrix, which has size $64\times64$ and a numerical rank of $10$. 
We test the diffusion completion method under principal submatrix masking with sampling rate $s=1\%$, and matrix completion method under both principal submatrix masking and random masking with various sampling rates, see Figure ~\ref{fig:comparison_completion_DM_MC}. 
We find that the diffusion model is able to reconstruct the off-diagonal block of DtN matrix within a relative error of $0.9\%$ with only $1\%$ of observed data.
In comparison, matrix completion fails to find any missing data under the same setting. This is because the masked measurement is already of low rank when the sampling rate $s=1\%$ is low, and thus becomes a solution to \Eqref{eq:mc_nucnorm}. Consequently, no missing data can be obtained from the nuclear norm optimization. 
We further test matrix completion under random masking with different sampling rate $s=1\%,15\%$ and $30\%$.
We observe that the same issue persists with low sampling rate $1\%$. An unsatisfactory reconstruction can be achieved with larger sampling rate $s=15\%$ with relative error $8.2\%$. 
When the random sampling rate is raised to $s=30\%$, a competitive reconstruction result with relative error $1.0\%$ is achieved. 
Overall, Figure \ref{fig:comparison_completion_DM_MC} illustrates that a learned distributional prior (diffusion) can interpolate across ultra-sparse observations, while low-rank matrix completion requires denser, randomly distributed observations to perform well. 

\begin{figure}[!htbp]
    \centering
    \includegraphics[width=1\linewidth]{data_completion/comparison_DM_MC.png}
    \caption{
    Diffusion completion versus matrix completion on an off-diagonal block under different masking patterns and sampling rate. \emph{Note that the matrix-completion baseline is only applied to off-diagonal blocks, where low-rank structure is expected, whereas our diffusion model is applied to the whole matrix and presents the off-block in this plot}
    \textbf{Left}: ground truth (normalized) upper-right block of DtN matrix. 
    \textbf{Right, top row}: masking matrices with sampling rate $s=1\%$, $s=15\%$ and $s=30\%$.
    \textbf{Right, bottom row}: corresponding reconstructions under the masks shown above. The diffusion model (first column) achieves high-quality reconstruction under the submatrix masking even at an extremely low sampling rate \(s=1\%\) with relative error (RE) \(0.9\%\). In contrast, matrix completion fails at \(s=1\%\) under both submatrix and random masking (second and third columns). Matrix completion requires substantially higher random sampling rates (fourth and fifth columns) to achieve comparable accuracy, reaching RE \(8.2\%\) at \(s=15\%\) and RE \(1.0\%\) at \(s=30\%\). In contrast, our diffusion model performs completion over the entire DtN matrix.
    }
    \label{fig:comparison_completion_DM_MC}
\end{figure}



\subsection{Results for inverse problem solving}\label{sec:result_IP}

In this section, we investigate the effectiveness of inverse problem solving with completed measurement data from the proposed method in Sections~\ref{sec:random_disk} and~\ref{sec:SL}, respectively. 
Given sparse DtN measurements, we evaluate the reconstruction performance of different methods, including \textbf{i)} a direct reconstruction by a sparse-data solver; \textbf{ii)} a level-$3$ hierarchical matrix completion followed by a full-data solver; and \textbf{iii)} a diffusion completion followed by the same full-data solver. 
Note that even with the ground truth data, the inherent nonlinearity and ill-posedness of the EIT problem can lead to non-negligible reconstruction errors, particularly under noisy or sparse measurement conditions. 
To this end, the output of a full-data solver with full data represents the best possible reconstruction from a specific inverse problem solver. We report this reconstruction as a reference to demonstrate the effectiveness of the aforementioned methods.

Specifically, to obtain reasonable reconstructions by matrix completion, we consider a higher sampling rate ($15\%$) in each off-diagonal block. This results in a global hierarchical random mask with an overall sampling rate of $25.625\%$. 
For diffusion completion and direct sparse solver, a more practical principal submatrix masking is considered with sampling rate $1\%$. 
These two different masking matrices are plotted in Figure \ref{fig:mask_comparison}.
\begin{figure}[!htbp]
    \centering
    \includegraphics[width=0.7\linewidth]{data_completion/mask_comparison.png}
\caption{Global masking matrices. \textbf{Left:} The mask matrix used for diffusion completion and the sparse-data solver. This mask consists of nonzero values over a principal submatrix with sampling rate $s=1\%$. \textbf{Right:} The hierarchical masking matrix with level $l=3$. In this case, all diagonal blocks are fully observed, and off-diagonal blocks are randomly sampled with sampling rate $s=15\%$, yielding an overall sampling rate of $25.625\%$.}

    \label{fig:mask_comparison}
\end{figure}


\subsubsection{Random disks}\label{sec:random_disk}
We present a visual comparison of reconstructions over different samples in Figure~\ref{fig:disk_deepsolver_noiseless} with clean data, and Figure~\ref{fig:disk_deepsolver_noise5e-2} with $5\%$ noise. A summary of errors over the testing dataset in both cases is shown in Table~\ref{tab:disk_recon_combined}.
In the noiseless setting, the full-data solver with full measurements achieves the best accuracy across all metrics, representing the solver’s best attainable performance. 
Using the same solver with diffusion-completed measurements from sparse data yields the closest performance to the best possible reconstructions. 
The direct sparse-data solver lags behind across all metrics, consistent with the observed blur and boundary erosion. Although H-matrix completion outperforms the direct sparse-data solver and is slightly worse than diffusion completion, it requires a much higher sampling rate (\(25.625\%\)) compared to the other cases (\(1\%\)). 
With \(5\%\) noise, all methods degrade slightly; diffusion completion remains superior to both the direct sparse-data solver and H-matrix completion.

\begin{table}[!htbp]
    \centering
    \caption{Reconstruction results for the \textbf{Disks} dataset. 
    \textbf{First row:} best possible reconstruction using full data.
    \textbf{Last three rows:} various inverse problem solving strategies with sparse data. The overall sampling rate is indicated in the brackets, and the best results with sparse data are marked in bold. }
    \label{tab:disk_recon_combined}
    \renewcommand{\arraystretch}{1.25}
    \begin{tabular}{@{}lcccccc@{}}
        \toprule
        \multirow{2}{*}{Method} &
        \multicolumn{3}{c}{Noiseless DtN} &
        \multicolumn{3}{c}{5\% Noisy DtN} \\
        \cmidrule(lr){2-4}\cmidrule(lr){5-7}
        & SSIM$\uparrow$ & RE$\downarrow$  & MAE$\downarrow$  
        & SSIM$\uparrow$  & RE$\downarrow$   & MAE$\downarrow$  \\
        \midrule
        full-data solver ($100\%$)        & 0.819 & 20.7\% & 0.176  & 0.808 & 21.6\% & 0.186  \\
        \hline
        diffusion completion ($1\%$) + full-data solver & \textbf{0.774} & \textbf{33.4}\% & \textbf{0.288}  & \textbf{0.762} & \textbf{34.4}\% & \textbf{0.302}  \\
 
        direct sparse-data solver ($1\%$)                     & 0.597 & 44.9\% & 0.495 & 0.578 & 46.5\% & 0.524 \\
        H-matrix completion ($25.625\%$) + full-data solver & 0.730 & 41.4\% & 0.361 & 0.712& 43.7\% & 0.388\\
        
    
        \bottomrule
    \end{tabular}
\end{table}

\begin{figure}[!htbp]
    \centering
    \includegraphics[width=\linewidth]{random_disks/random_disk_all_noiseless.png}

\caption{
Reconstructions from the \emph{noiseless} \textbf{random disks} dataset. 
\textbf{Columns (left$\to$right):} (1) \emph{Ground truth} conductivity; (2) \emph{Best possible reconstruction} from the full-data solver (overall sampling rate $100\%$); (3) Reconstructions from the full-data solver applied to \emph{diffusion completion} of measurements (overall sampling rate $1\%$); (4) Reconstructions from the \emph{direct sparse-data solver} (overall sampling rate $1\%$); (5) Reconstructions from the full-data solver applied to \emph{H-matrix completion} of measurements (overall sampling rate $25.625\%$). \textbf{Rows (top$\to$bottom):} Examples with number of disks $n=2,3,4,5$.}
\label{fig:disk_deepsolver_noiseless}
\end{figure}

\begin{figure}[!htbp]
    \centering
    \includegraphics[width=\linewidth]{random_disks/random_disk_all_noise5e-2.png}
\caption{
Reconstructions from the \textbf{random disks} dataset with \emph{$5\%$ noise}. 
\textbf{Columns (left$\to$right):} (1) \emph{Ground truth} conductivity; (2) \emph{Best possible reconstruction} from the full-data solver (overall sampling rate $100\%$); (3) Reconstructions from the full-data solver applied to \emph{diffusion completion} of measurements (overall sampling rate $1\%$); (4) Reconstructions from the \emph{direct sparse-data solver} (overall sampling rate $1\%$); (5) Reconstructions from the full-data solver applied to \emph{H-matrix completion} of measurements (overall sampling rate $25.625\%$). \textbf{Rows (top$\to$bottom):} Examples with number of disks $n=2,3,4,5$.}



\label{fig:disk_deepsolver_noise5e-2}
\end{figure}


\subsubsection{Random Shepp-Logan phantoms}\label{sec:SL}
The same experiments and comparisons are conducted for the Shepp--Logan dataset with visual reconstructions in Figs.~\ref{fig:sl_deepsolver_noiseless} and \ref{fig:sl_deepsolver_noise5e-2}, and errors shown in Table~\ref{tab:sl_recon_combined}. In the noiseless setting, diffusion completion + full-data solver produces reconstructions close to the best possible result, preserving the thin skull boundary and the intracranial cavities. Unlike the Disks dataset, the {direct sparse-data solver} performs better than {H-matrix completion} on Shepp--Logan: the {direct sparse-data solver} recovers the main structures but tends to smooth thin features, whereas {H-matrix completion + full-data solver} fails to recover both the skull boundary and the cavities. With \(5\%\) noise, robustness differences become more apparent. The {best possible reconstruction} and the {diffusion-completion} strategy both remain stable, while the {direct sparse-data solver} exhibits edge distortions at the boundary and amplified blur artifacts in the zoomed region. The {H-matrix completion + full-data solver} continues to fail as in the noiseless setting. Overall, diffusion completion narrows the gap to the best possible reconstruction across all metrics and maintains this advantage under noise, whereas the direct sparse-data solver and H-matrix completion struggle to solve the inverse problem.

\begin{table}[!htbp]
    \centering
    \caption{Reconstruction results for the \textbf{Shepp--Logan} dataset. 
    \textbf{First row:} best possible reconstruction using full data.
    \textbf{Last three rows:} various inverse problem solving strategies with sparse data. The overall sampling rate is indicated in the brackets, and the best results with sparse data are marked in bold.}
    \label{tab:sl_recon_combined}
    \renewcommand{\arraystretch}{1.25}
    \begin{tabular}{@{}lcccccc@{}}
        \toprule
        \multirow{2}{*}{Method} &
        \multicolumn{3}{c}{Noiseless DtN} &
        \multicolumn{3}{c}{5\% Noisy DtN} \\
        \cmidrule(lr){2-4}\cmidrule(lr){5-7}
        & SSIM$\uparrow$ & RE$\downarrow$ & MAE$\downarrow$ 
        & SSIM$\uparrow$ & RE$\downarrow$ & MAE$\downarrow$ \\
        \midrule
        full-data solver ($100\%$)                 
        & 0.930 & 8.2\% & 0.030 & 0.898 &10.2\% &0.042 \\
        \hline
        diffusion completion ($1\%$) + full-data solver & \textbf{0.910} & \textbf{10.3}\% & \textbf{0.036} & \textbf{0.842} & \textbf{14.3\%} & \textbf{0.055} \\

        direct sparse-data solver ($1\%$)                     
        &0.886 & 11.3\% & 0.043 & 0.776 & 16.4\% & 0.075 \\
        H-matrix completion ($25.625\%$) + full-data solver & 0.825 & 17.0\% & 0.067 & 0.788 & 19.5\% & 0.085\\
        \bottomrule
    \end{tabular}
\end{table}

\begin{figure}[!htbp]
    \centering
    \includegraphics[width=\linewidth]{shepp_logan/shepp_logan_noiseless_conductivity.png}
\caption{
Reconstructions from the \emph{noiseless} \textbf{Shepp--Logan} dataset. 
\textbf{Columns (left$\to$right):} (1) \emph{Ground truth} conductivity; (2) \emph{Best possible reconstruction} from the full-data solver (overall sampling rate $100\%$); (3) Reconstructions from the full-data solver applied to \emph{diffusion completion} of measurements (overall sampling rate $1\%$); (4) Reconstructions from the \emph{direct sparse-data solver} (overall sampling rate $1\%$); (5) Reconstructions from the full-data solver applied to \emph{H-matrix completion} of measurements (overall sampling rate $25.625\%$). \textbf{Rows (top$\to$bottom):} full field of view (top) with a red box indicating the zoom region; zoomed-in crops of that region (bottom).
}

\label{fig:sl_deepsolver_noiseless}
\end{figure}

\begin{figure}[!htbp]
    \centering
    \includegraphics[width=\linewidth]{shepp_logan/shepp_logan_noise5e-2_conductivity.png}
\caption{
Reconstructions from the random \textbf{Shepp--Logan} dataset \emph{$5\%$ noise}. 
\textbf{Columns (left$\to$right):} (1) \emph{Ground truth} conductivity; (2) \emph{Best possible reconstruction} from the full-data solver (overall sampling rate $100\%$); (3) Reconstructions from the full-data solver applied to \emph{diffusion completion} of measurements (overall sampling rate $1\%$); (4) Reconstructions from the \emph{direct sparse-data solver} (overall sampling rate $1\%$); (5) Reconstructions from the full-data solver applied to \emph{H-matrix completion} of measurements (overall sampling rate $25.625\%$). \textbf{Rows (top$\to$bottom):} full field of view (top) with a red box indicating the zoom region; zoomed-in crops of that region (bottom).
}
    \label{fig:sl_deepsolver_noise5e-2}
\end{figure}


\section{Conclusion and Discussion}

We have proposed a diffusion-based DtN completion framework for EIT with extremely sparse boundary measurements and evaluated it on the Disks and Shepp--Logan datasets. Using the same inverse solver across all settings, reconstructions from diffusion completion are consistently closest to the best possible reconstruction that uses full measurements among the direct sparse-data solver and H-matrix completion. These results support our diffusion completion design that improves stability and accuracy without modifying the inverse solvers on-the-shelf.
Our analysis shows that the discretized DtN measurement set $F_N(K)$ inherits the low-dimensional
structure of the polygon conductivity class $K$. The intrinsic complexity of DtN measurements can be controlled via a
metric entropy bound determined by the underlying parameter dimension $2n_{\mathrm{v}}$, rather than by the ambient measurement dimension. Combining this complexity estimate with the conditional diffusion-model convergence theorem yields a
nonasymptotic end-to-end completion guarantee in our EIT setting, with an explicit dependence on the diffusion steps $T$ and vertices $n_{\mathrm{v}}$.

We next outline several promising directions for further development. First, although we considered simulated data in this paper, diffusion completion can be directly applied to real EIT application provided sufficient lab data are available for training. Second, encoding structure of the DtN matrix can improve generalization and data efficiency. In addition to hierarchical low-rankness \citep{bebendorf2003existence, ballani2017matrices}, DtN measurements in our setting are symmetric and satisfy zero row or column sum constraints \citep{somersalo1992existence}, which can be enforced by specific network architectural designs or training penalties. Third, accelerating sampling is important for practical deployment. Distillation \citep{song2023consistency} and few-step samplers \citep{song2020denoising} in diffusion models can be directly incorporated into our framework and reduce inference cost while maintaining accuracy on conditional tasks. Fourth, the generative measurement prior can guide the choice of boundary signals. Expected information gain provides a clear objective for selecting informative measurements, and modern Bayesian optimal experimental design methods make this feasible in EIT settings \citep{dong2025variational, hellmuth2025data}.



\section{Acknowledgements}
The authors were partially supported by the US National Science Foundation under awards DMS-2244988, CAIG RISE-2530596, IIS-2520978, and the DARPA D24AP00325-00.
   
\newpage
\bibliographystyle{plainnat}
\bibliography{ref}  

\newpage
\appendix
\section{Uniform Lipschitz boundary for convex admissible polygons}
\label{sec:lipschitz_bdy_class}
\begin{definition}[Lipschitz boundary class]
\label{def: lipschitz_bdy_class}
    Fix constants $r_0>0$ and $K_0> 1$.
We say that an open set $D\subset\R^2$ has \emph{Lipschitz boundary of class $(r_0,K_0)$} if for every point
$x_0\in\partial D$ there exist a rigid motion $\mathcal T_{x_0}$ (a translation followed by a rotation) with
$\mathcal T_{x_0}(x_0)=0$ and a $K_0$-Lipschitz function $\phi:\,[-r_0,r_0]\to\R$ such that
\[
\mathcal T_{x_0}(D)\cap\big([-r_0,r_0]\times[-K_0r_0,K_0r_0]\big)
=\{(x_1,x_2): x_1\in[-r_0,r_0],\ x_2>\phi(x_1)\}\,.
\]
\end{definition}

\begin{lemma}[Uniform Lipschitz boundary for convex admissible polygons]\label{lem:A4_from_A23_conv}
Let $P=P(v) \in \mathcal{A}$, $\partial P$ is of Lipschitz class $(r_0,K_0)$ with the \emph{uniform} choices
\[
r_0:=\frac{d_1}{4},
\qquad
K_0:=\cot\!\Big(\frac{\beta_0}{2}\Big).
\]
\end{lemma}

\begin{proof}
Fix $x_0\in\partial P$. We distinguish two cases.

\emph{Case 1: $x_0$ lies in the relative interior of an edge.}
Let $e=[v_i,v_{i+1}]$ be that edge. Since $|v_{i+1}-v_i|\ge d_1$, the ball
$B_2(x_0,r_0)$ with $r_0=d_1/4$ intersects $\partial P$ only along the segment of $e$
(no vertex is contained in the ball).
After a rigid motion sending $x_0$ to $0$ and rotating so that $e$ becomes $\{x_2=0\}$,
the convexity implies that in a small neighborhood $P$ coincides with either $\{x_2>0\}$
or $\{x_2<0\}$, hence the Lipschitz representation holds with Lipschitz constant $0$.

\emph{Case 2: $x_0$ is a vertex, say $x_0=v_i$.}
Translate $v_i$ to $0$. Let the two incident edges be the rays along the lines
$\ell_-$ and $\ell_+$ meeting at interior angle $\beta_i\in[\beta_0,\pi)$.
Rotate coordinates so that the bisector of the interior angle is the $x_2$-axis.
Then each line $\ell_\pm$ makes an angle $\beta_i/2$ with the $x_2$-axis, hence an angle
$\pi/2-\beta_i/2$ with the $x_1$-axis. Therefore each $\ell_\pm$ is the graph of an affine function
$x_2=\phi_\pm(x_1)$ with slope bounded by
\[
|\phi_\pm'(x_1)|=\tan\!\Big(\frac{\pi}{2}-\frac{\beta_i}{2}\Big)=\cot\!\Big(\frac{\beta_i}{2}\Big)
\le \cot\!\Big(\frac{\beta_0}{2}\Big)=K_0.
\]
Since $|v_{i\pm 1}-v_i|\ge d_1$, the ball $B_2(0,r_0)$ with $r_0=d_1/4$ intersects $\partial P$
only along these two edges. In the above coordinates, $P\cap B_2(0,r_0)$ coincides with the epigraph
$\{x_2>\phi(x_1)\}$ where $\phi:=\max\{\phi_-,\phi_+\}$, and $\phi$ is $K_0$-Lipschitz.

Combining the two cases gives that $\partial P$ is of Lipschitz class $(r_0,K_0)$ with the stated
uniform constants.
\end{proof}

\section{Proof for Lemma~\ref{lem:rho_s}}
\label{sec:proof_rho_s}
We first show a useful fact that perturbing vertices by $\delta$ leads to at most $\delta$ perturbation around the polygonal boundary in Hausdorff distance.

\begin{lemma}[Boundary perturbation under vertex perturbations]\label{lem:dist-boundary-perturb}
Let $v,w\in \mathcal{V}^{n_{\mathrm v}}$ be counterclockwise ordered vertex, then
\[
d_H\big(\partial P(v),\partial P(w)\big)\ \le\ \|v-w\|_{\mathrm{poly}}.
\]
Moreover, if $P(v)\subset\Omega$ and $P(w)\subset\Omega$, then
$\dist(P(u),\partial\Omega)=\dist(\partial P(u),\partial\Omega)$ for $u=v,w$\,, and
\[
\dist\big(P(v),\partial\Omega\big)\ \ge\ \dist\big(P(w),\partial\Omega\big)\ -\ \|v-w\|_{\mathrm{poly}}.
\]
\end{lemma}

\begin{proof}
Fix $i$ and consider the segments $e_i(v)=[v_i,v_{i+1}]$ and $e_i(w)=[w_i,w_{i+1}]$.
For any $t\in[0,1]$, set $x(t)=(1-t)v_i+t v_{i+1}\in e_i(v)$ and $y(t)=(1-t)w_i+t w_{i+1}\in e_i(w)$.
Then
\[
|x(t)-y(t)|\le (1-t)|v_i-w_i|+t|v_{i+1}-w_{i+1}|\le \|v-w\|_{\mathrm{poly}}.
\]
Consequently, we have
\[
\dist\!\big(x(t),e_i(w)\big)\le \|v-w\|_{\mathrm{poly}}. 
\]
Taking the supremum over $t\in[0,1]$ yields
\begin{equation}\label{eqn:sup_x_e_i}
\sup_{x\in e_i(v)}\dist(x,e_i(w))\le \|v-w\|_{\mathrm{poly}}\,.
\end{equation}
Consider $x\in\partial P(v)$, then $x\in e_i(v)$ for some $i$, and since $e_i(w)\subset\partial P(w)$ we have
\[
\dist(x,\partial P(w))\le \dist(x,e_i(w))\le \|v-w\|_{\mathrm{poly}}\,,
\]
where the last inequality results from \Eqref{eqn:sup_x_e_i}.
Taking the supremum over $x\in\partial P(v)$ gives
$\sup_{x\in\partial P(v)}\dist(x,\partial P(w))\le \|v-w\|_{\mathrm{poly}}$.
Since $\partial P(v)$ and $\partial P(w)$ are unions of the corresponding edges, we also obtain
$\sup_{x\in\partial P(v)}\dist(x,\partial P(w))\le \|v-w\|_{\mathrm{poly}}$.
By symmetry (exchanging $v$ and $w$) we also have
$\sup_{y\in\partial P(w)}\dist(y,\partial P(v))\le \|v-w\|_{\mathrm{poly}}$, proving the Hausdorff estimate. 

For $P(v)\subset\Omega$, since $P(v)$ is compact and $\partial\Omega$ is closed, 
any point $x\in P(v)$ minimizing $\dist(x,\partial\Omega)$ must lie on $\partial P(v)$:
if $x$ were in the interior of $P(v)$, one could move slightly toward a nearest boundary point of $\partial\Omega$
while staying inside $P(v)$ and decrease the distance, a contradiction. Hence
$\dist(P(v),\partial\Omega)=\dist(\partial P(v),\partial\Omega)$.

For the distance inequality, note that if $d_H(A,A')\le \delta$ for nonempty closed sets, then
$\dist(A,B)\ge \dist(A',B)-\delta$ for any closed $B$, by the triangle inequality in $\R^2$. 
\end{proof}

Now we can prove Lemma ~\ref{lem:rho_s} as follows. 
\begin{proof}[Proof of Lemma ~\ref{lem:rho_s}]
Write $e_i^\ast\coloneqq [v_i^\ast,v_{i+1}^\ast]$ for the edges of $\partial P^\ast$.
Since $P^\ast$ is convex, $\partial P^\ast$ is simple, and every pair of nonadjacent edges is disjoint.
Hence
\[
\eta^\ast\coloneqq \tfrac12\min\big\{\dist(e_i^\ast,e_j^\ast):\ e_i^\ast,e_j^\ast\ \text{nonadjacent}\big\}>0.
\]
Moreover, since $P^\ast\in\mathcal A_{1/2}$, we have the strict margins
\[
m_0\coloneqq \dist(P^\ast,\partial\Omega)-\tfrac12 d_0>0,\qquad
m_1\coloneqq \min_i |v_{i+1}^\ast-v_i^\ast|-\tfrac12 d_1>0,
\qquad
\beta_s\coloneqq \min_i\Big\{\beta_i(v^\ast)-\tfrac12\beta_0,\ \pi-\tfrac12\beta_0-\beta_i(v^\ast)\Big\}>0.
\]

For any $i$,
\[
\big||v_{i+1}-v_i|-|v_{i+1}^\ast-v_i^\ast|\big|
\le |(v_{i+1}-v_{i+1}^\ast)-(v_i-v_i^\ast)|
\le 2\|v-v^\ast\|_{\mathrm{poly}}.
\]
Thus if $\|v-v^\ast\|_{\mathrm{poly}}<m_1/4$, then $|v_{i+1}-v_i|>d_1/2$ for all $i$, i.e.\ \textup{(A2$'$)} holds.


For each edge $e_i(v)=[v_i,v_{i+1}]$, the same interpolation argument as in Lemma~\ref{lem:dist-boundary-perturb}
gives $d_H(e_i(v),e_i^\ast)\le \|v-v^\ast\|_{\mathrm{poly}}$.
Hence for any nonadjacent edge index pair $(i,j)$,
\[
\dist(e_i(v),e_j(v))
\ \ge\ \dist(e_i^\ast,e_j^\ast)-d_H(e_i(v),e_i^\ast)-d_H(e_j(v),e_j^\ast)
\ \ge\ 2\eta^\ast-2\|v-v^\ast\|_{\mathrm{poly}}.
\]
If $\|v-v^\ast\|_{\mathrm{poly}}<\eta^\ast$, then $\dist(e_i(v),e_j(v))>0$ for all nonadjacent pairs,
so nonadjacent edges cannot intersect. Moreover, once we also enforce the angle lower bound $\beta_i(v)>\beta_0/2$ for all $i$ (proved below),
adjacent edges cannot overlap. Hence $\partial P(v)$ is a simple closed polygonal curve.


Now we show the inclusion in $\Omega$ and the boundary margin.
Since $P^\ast\subset\Omega$ and $\dist(P^\ast,\partial\Omega)>d_0/2$, every vertex satisfies
$\dist(v_i^\ast,\partial\Omega)\ge \dist(P^\ast,\partial\Omega)>0$.
Set
\[
\eta_\Omega^\ast\coloneqq \tfrac12\,\dist(P^\ast,\partial\Omega)>0.
\]
If $\|v-v^\ast\|_{\mathrm{poly}}<\eta_\Omega^\ast$, then $|v_i-v_i^\ast|<\eta_\Omega^\ast$ for all $i$,
hence $v_i\in\Omega$ for all $i$. As $\Omega$ is convex, each edge $[v_i,v_{i+1}]\subset\Omega$ and therefore
$P(v)\subset\Omega$. Furthermore, by Lemma~\ref{lem:dist-boundary-perturb},
\[
\dist(P(v),\partial\Omega)\ge \dist(P^\ast,\partial\Omega)-\|v-v^\ast\|_{\mathrm{poly}}.
\]
Thus if $\|v-v^\ast\|_{\mathrm{poly}}<m_0$, then $\dist(P(v),\partial\Omega)>d_0/2$, i.e.\ \textup{(A1$'$)} holds.

We now consider the angle bounds and convexity in (A3$'$). On the set where \textup{(A2$'$)} holds, the interior angles $\beta_i(v)$ are continuous functions of $v$
(e.g.\ they can be expressed via dot/cross products of the adjacent edge vectors).
Hence there exists $\rho_{\mathrm{ang}}>0$ such that if $\|v-v^\ast\|_{\mathrm{poly}}<\rho_{\mathrm{ang}}$ then
$|\beta_i(v)-\beta_i(v^\ast)|<\beta_s/2$ for all $i$, and consequently
$\beta_0/2<\beta_i(v)<\pi-\beta_0/2$ for all $i$, which yields \textup{(A3$'$)}.

Finally, define
\[
\rho_{P^\ast}\coloneqq \min\Big\{\eta^\ast,\ \eta_\Omega^\ast,\ m_0,\ m_1/4,\ \rho_{\mathrm{ang}}\Big\}.
\]
Then $\|v-v^\ast\|_{\mathrm{poly}}<\rho_{P^\ast}$ implies \textup{(A1$'$)--(A3$'$)}.
\end{proof}

\section{Proof of Lemma~\ref{lem:chart-M}}
\label{sec:chart-M}

\begin{lemma}\label{lem:L1_to_H_poly}
There exists a constant $C_H>0$, depending only on the a priori data, such that for all
$P,Q\in\mathcal A_{1/2}$,
\[
d_H(\partial P,\partial Q)\ \le\ C_H\|\gamma_P-\gamma_Q\|_{L^1(\Omega)}^{1/2}.
\]
\end{lemma}

\begin{proof}[Proof of Lemma ~\ref{lem:L1_to_H_poly}]
Let $\delta:=d_H(\partial P,\partial Q)$. It suffices to prove the case when $\delta>0$. Without loss of generality, we can choose $x\in\partial P$ such that $\dist(x,\partial Q)=\delta$ since $\partial P$ is compact. Consider open ball $\tilde{B}_2(x,\delta)$, $\tilde{B}_2(x,\delta)\cap\partial Q=\varnothing$, hence $\chi_Q$ is constant on $\tilde{B}_2(x,\delta)$.

Since $\partial P$ is Lipschitz of class $(r_0,K_0)$, at $x$ there exist local coordinates in which
$P$ coincides with $\{(x_1,x_2): x_2>\phi(x_1)\}$, with $\phi$ $K_0$-Lipschitz.
In these coordinates, for any $h\le \min\{r_0,\delta/2\}$ the truncated cone
\[
C^+(h):=\{(x_1,x_2): 0<x_2<h,\ |x_1|<x_2/K_0\}
\]
is contained in $P\cap \tilde{B}_2(x,\delta)$, and the reflected cone
\[
C^-(h):=\{(x_1,x_2): -h<x_2<0,\ |x_1|<-x_2/K_0\}
\]
is contained in $(\mathbb R^2\setminus P)\cap \tilde{B}_2(x,\delta)$.
Moreover $|C^+(h)|=|C^-(h)|=h^2/K_0$.

If $\chi_Q\equiv 0$ on $\tilde{B}_2(x,\delta)$ then $C^+(h)\subset P\setminus Q\subset P\Delta Q$; if $\chi_Q\equiv 1$ on $\tilde{B}_2(x,\delta)$ then $C^-(h)\subset Q\setminus P\subset P\Delta Q$.
In both cases, $|P\Delta Q|\ge h^2/K_0$ for $h=\min\{r_0,\delta/2\}$, which implies
$\delta\le 2\sqrt{K_0}\,|P\Delta Q|^{1/2}$ whenever $\delta\le 2r_0$.
If $\delta>2r_0$, the same estimate with $h=r_0$ gives $|P\Delta Q|\ge r_0^2/K_0$ and thus
$\delta\le \diam(\Omega)\le \diam(\Omega)\sqrt{K_0}/r_0\cdot |P\Delta Q|^{1/2}$.
Choosing $C_H = (\kappa-1)^{-1/2}\max \{2\sqrt{K_0}, \diam(\Omega)\sqrt{K_0}/r_0\}$ yields this lemma. 
\end{proof}


\begin{lemma}\label{lem:vertex_stability_poly}
There exist $C_{\mathrm{vs}}>0$ depending only on the a priori data in Remark~\ref{rem:apriori_poly} such that, if $P^0,P^1\in\mathcal A_{1/2}$ with corresponding vertex list $v^0$ and $v^1$, then their vertices can be ordered in such a way that 
\[
\|v^0-v^1\|_{\mathrm{poly}} \le\ C_{\mathrm{vs}}\, d_H(\partial P^0,\partial P^1).
\]
\end{lemma}
\begin{proof}[Proof of Lemma ~\ref{lem:vertex_stability_poly}]
Let $\widetilde{\mathcal{A}}_{1/2}$ be the non-strict relaxed class with parameters $(d_0/2, d_1/2, \beta_0/2)$, i.e., $\widetilde{\mathcal{A}}_{1/2} = \{P\subset \Omega, \dist(P, \partial \Omega) \geq d_0/2, |v_{i+1}-v_{i}|\geq d_1/2, \frac{\beta_0}{2} \leq \beta_i(v)\leq \pi-\frac{\beta_0}{2}, \partial P \text{ is simple} \}$. Therefore $\mathcal{A}_{1/2} \subset \widetilde{\mathcal{A}}_{1/2}$. 
By Proposition 3.3 in \citet{beretta2022global}, there exist constants $\delta_0>0$ and $C_{\mathrm{0}}>0$,
depending only on the a priori data such that, if $P^0,P^1\in \widetilde{\mathcal{A}}_{1/2} $ satisfy $d_H(\partial P^0,\partial P^1)\le \delta_0$, then the vertices can be ordered in such a way that 
\[
\|v^0-v^1\|_{\mathrm{poly}} \leq C_{0}\, d_H(\partial P^0,\partial P^1).
\]
It suffices to prove the case when $d_H(\partial P^0,\partial P^1) > \delta_0$. Suppose $d_H(\partial P^0,\partial P^1) > \delta_0$, 
\[
\|v^0-v^1\|_{\mathrm{poly}} \leq \|v^0\|_{\mathrm{poly}}+\|v^1\|_{\mathrm{poly}} \leq 2 \diam(\Omega) \frac{d_H(\partial P^0,\partial P^1)}{\delta_0} \,.
\]
Taking $C_{\mathrm{vs}} = \max(C_0,\frac{2\diam(\Omega)}{\delta_0} )$ finishes the proof.
\end{proof}

\begin{lemma}[$L^1$-stability under vertex perturbations]\label{lem:symdiff_poly}
Fix $P^\star\in\mathcal A_{1/2}$ and consider the neighborhood $\varphi_{P^\star}(U_{P^\star})$
from Lemma~\ref{lem:chart-M}. Let $v,w\in\varphi_{P^\star}(U_{P^\star})$ and set
$\delta\coloneqq \|v-w\|_{\mathrm{poly}}$.
Then
\[
\|\chi_{P(v)}-\chi_{P(w)}\|_{L^1(\Omega)}
= |P(v)\Delta P(w)| \le C\,\delta\,.
\]
Here $C=2n_{\mathrm v}(\pi+1) \diam(\Omega)>0$ is a constant.
\end{lemma}
\begin{proof}[Proof of Lemma~\ref{lem:symdiff_poly}]
Consider the segment $v(t)=(1-t)v+tw$ for $t\in[0,1]$ and write $P_t\coloneqq P(v(t))$. 
By Lemma~\ref{lem:chart-M}, all $v(t)$ are relaxed admissible, thus $P_t$ is well-defined and $P_t\subset\Omega$.
Let $x\in P(v)\Delta P(w)$, we first show there exists $t^\ast\in[0,1]$ such that $x\in\partial P_{t^\ast}$.
Define
\[
g(t)\coloneqq \dist(x,P_t^c)-\dist(x,P_t),
\]
where the map $t\mapsto g(t)$ is continuous.
If $g(0)=0$ or $g(1)=0$, $x \in \partial P(v)$ or $x \in \partial P(w)$; otherwise  $g(0)$ and $g(1)$ have opposite signs because $x$ belongs
to exactly one of $P_0$ and $P_1$. Hence there exists $t^\ast\in(0,1)$ such that $g(t^\ast)=0$, which implies
$\dist(x,P_{t^\ast})=\dist(x,P_{t^\ast}^c)=0$ and therefore $x\in\partial P_{t^\ast}$.

By Lemma~\ref{lem:dist-boundary-perturb} we have
\[
d_H(\partial P_{t^\ast},\partial P(v))
\le \|v(t^\ast)-v\|_{\mathrm{poly}}\le \delta,
\]
so $\dist(x,\partial P(v))\le\delta$. This shows
\[
P(v)\Delta P(w)\subset \{x\in\R^2:\dist(x,\partial P(v))\le\delta\}.
\]

Now $\partial P(v)$ is a union of $n_{\mathrm v}$ line segments. The $\delta$-neighborhood of a segment of
length $\ell$ has area at most $2\delta\,\ell+\pi\delta^2$ (rectangle plus two half-disks).
Summing over all edges yields
\[
|P(v)\Delta P(w)|
\le 2\delta\,\Per(P(v)) + n_{\mathrm v}\pi\delta^2.
\]
Finally, $\Per(P(v))\le n_{\mathrm v}\diam(\Omega)$ and $\delta\le 2R_{P^\star}\leq 2 \diam(\Omega)$ on the chart, hence $|P(v)\Delta P(w)| \le C\,\delta$ with $C=2n_{\mathrm v}(\pi+1) \diam(\Omega)$.
\end{proof}

We are now ready to prove Lemma ~\ref{lem:chart-M}. 
\begin{proof}[Proof of Lemma ~\ref{lem:chart-M}]

If $\|v-v^\star\|_{\mathrm{poly}}<R_{P^\star}$, then $R_{P^\star}\le \rho_{P^\star}$ implies $P(v)\in\mathcal A_{1/2}$
by Lemma~\ref{lem:rho_s}, hence $\gamma_{P(v)}\in M$ and $U_{P^\star}\subset M$. 

Moreover, $R_{P^\star}\le \frac14\min_{i\neq j}|v_i^\star-v_j^\star|$ implies that the balls
$B_2(v_i^\star,R_{P^\star})$ are pairwise disjoint. Therefore each perturbed vertex $v_i$ lies in exactly one ball of $v_i^\star$,
so no cyclic shift or reversal of the vertex labeling is possible inside $U_{P^\star}$; thus $\varphi_{P^\star}$ is well-defined
injective.
Next we show that $\varphi_{P^\star}$ is a homeomorphism.

(continuity of $\varphi_{P^\star}^{-1}$).
By Lemma~\ref{lem:symdiff_poly}, the map
\[
\varphi_{P^\star}^{-1}:\mathcal O_{P^\star}\to L^1(\Omega),\qquad v\mapsto \gamma_{P(v)},
\]
is (locally) Lipschitz in $L^1(\Omega)$ on $\mathcal O_{P^\star}$, hence continuous.

(continuity of $\varphi_{P^\star}$).
For each $\gamma_{P_0}$  in $U_{P^\star}$, for any sequence $\gamma_{P_k}$ in $U_{P^\star}$ converges to $\gamma_{P_0}$, $|P_k \Delta P_0| = \frac{1}{(\kappa-1)}\|\gamma_{P_k}-\gamma_{P_0} \|_{L^1(\Omega)} \to 0$. By Lemma~\ref{lem:L1_to_H_poly}, $d_{H}(\partial P_k, \partial P_0)\to 0$, and then by Lemma~\ref{lem:vertex_stability_poly}, we have $\|v^{k} - v^0\|_{\mathrm{poly}} \to 0$ for some $v^k$ and $v^0$ such that $P(v^k)=P_k$ and $P(v^0)=P_0$. Thus  $\varphi_{P^\star}$ is continous. 

($U_{P^\star}$ is open in $M$).
For each $\gamma_{P_0}$  in $U_{P^\star}$, we show that $\gamma_{P_0}$ is the interior point in $U_{P^\star}$. Assume by contradiction that for every $k>=1$, there exists $\gamma_{P_k} \in M \cap B_{L^1}(\gamma_{P_0}, \frac{1}{k})$ such that  $\gamma_{P_k} \notin U_{P^\star}$. Then $\|v^k-v^\star\|_{\mathrm{poly}}\geq R_{P^\star}$ for some $v^k$ such that $P(v^k) = P_k$. By the triangle inequality we have
\begin{equation}
\label{eq:contract_U_open}
    0 < R_{P^\star} - \|v^0-v^\star\|_{\mathrm{poly}} \leq \|v^k-v^0\|_{\mathrm{poly}}, \qquad k\geq 1 \,,
\end{equation}


for some $v^0$ such that $P(v^0) = P_0$. Since $\|\gamma_{P_k}-\gamma_{P_0} \|_{L^1(\Omega)} \to 0$, we have $\|v^k-v^0\|_{\mathrm{poly}} \to 0$, which contracts to \Eqref{eq:contract_U_open}.
\end{proof}






\section{Proof for Lemma~\ref{lem:K_compact_poly}}
\label{sec:K_compact_poly}
The following proof follows the idea in the Proof of theorem 3.1 from \cite{alberti2022inverse}.
\begin{proof}
Let $(\gamma_{P_j})_j\subset K$ with $P_j=P(v^j)$ and $v^j=(v^j_1,\dots,v^j_{n_{\mathrm v}})$
the counterclockwise vertex list.
Set $\Omega_{d_0}:=\{x\in\R^2:\dist(x,\partial\Omega)\ge d_0\}=\overline{B(0,1-d_0)}$.
By (A1) we have $P_j\subset \Omega_{d_0}$, hence $v^j_i\in\Omega_{d_0}$ for all $i,j$.

Since $\Omega_{d_0}$ is compact and there are only finitely many cyclic relabelings,
we may pass to a subsequence indexed by $j_m$ and relabel the vertices by a fixed cyclic shift so that $v^{j_m}\to v^\ast$ in $\R^{2n_{\mathrm v}}$.
Define $P^\ast:=P(v^\ast)$ and we show that $P^\ast\in\mathcal A$.
\begin{itemize}
\item[(A1)] Since $\Omega_{d_0}$ is convex, $P^\ast=\mathrm{conv}\{v^\ast_i\}\subset \Omega_{d_0}\subset\Omega$,
and therefore $\dist(P^\ast,\partial\Omega)\ge d_0$.
\item[(A2)] For each $i$, the map $v\mapsto |v_{i+1}-v_i|$ is continuous, hence
$|v^\ast_{i+1}-v^\ast_i|=\lim_j |v^j_{i+1}-v^j_i|\ge d_1$.
\item[(A3)] For each $i$, the interior angle $\beta_i(v)$ is a continuous function of the adjacent edge vectors
(as long as $|v_{i}-v_{i-1}|,|v_{i+1}-v_i|>0$, ensured by (A2)).
Thus $\beta_i(v^\ast)=\lim_j\beta_i(v^j)$ and
$\beta_0\le \beta_i(v^\ast)\le \pi-\beta_0$.
Moreover, the uniform bound $\beta_i(v^\ast)\le\pi-\beta_0$ rules out collinearity of adjacent edges,
so $P^\ast$ is a convex polygon with the same counterclockwise ordering, hence simple.
\end{itemize}
Thus $P^\ast\in\mathcal A$ and $\gamma_{P^\ast}\in K$.

Finally, by Lemma~\ref{lem:symdiff_poly}, for large $j_m$ we have
$|P_{j_m}\Delta P^\ast|\to 0$, hence
\[
\|\gamma_{P_{j_m}}-\gamma_{P^\ast}\|_{L^1(\Omega)}
=|\kappa-1|\,|P_j\Delta P^\ast|\to 0.
\]
Therefore every sequence in $K$ admits an $L^1(\Omega)$-convergent subsequence with limit in $K$,
so $K$ is compact in $L^1(\Omega)$.
\end{proof}


\section{Proof for Lemma~\ref{lem:finite-atlas}}
\begin{proof}[Proof of Lemma ~\ref{lem:finite-atlas}]
Since $M\subset L^{1}(\Omega)$ is endowed with the $L^{1}$-subspace topology (equivalently, the inclusion
$\iota:M\hookrightarrow L^{1}(\Omega)$ is a continuous embedding), compactness of $K$ in $L^{1}(\Omega)$ implies
that $K$ is compact in $M$.


For each $\gamma_P\in K$, the set
\[
V_{P}=\Big\{\gamma\in U_P:\ \|\varphi_P(\gamma)-v^{P}\|_{\mathrm{poly}}<\tfrac12 R_{P}\Big\}
=\varphi_P^{-1}\!\bigl(\{v:\|v-v^{P}\|_{\mathrm{poly}}<\tfrac12R_P\}\bigr)
\]
is open in $M$ (it is the preimage of an open ball under the continuous chart map $\varphi_P$).
Moreover $\gamma_P\in V_P$, hence $\{V_P\}_{\gamma_P\in K}$ is an open cover of $K$.
By compactness of $K$ in $M$, there exist $P_1,\dots,P_m\in\mathcal A$ such that
$K\subset\bigcup_{i=1}^m V_{P_i}$.

Fix $i\in\{1,\dots,m\}$ and by the choice of $R_{P_i}$ we have $B_i\subset \varphi_{P_i}(U_{P_i})$.
Since $B_i$ is compact in $\R^{2n_{\mathrm v}}$ and $\varphi_{P_i}^{-1}$ is continuous on $\varphi_{P_i}(U_{P_i})$,
the set $\varphi_{P_i}^{-1}(B_i)$ is compact (hence closed) in $M$ and contains $V_{P_i}$.
Therefore
\[
\overline{V}_{P_i}\subset \varphi_{P_i}^{-1}(B_i)\subset U_{P_i},
\]
so $\varphi_{P_i}$ is well-defined on $\overline{V}_{P_i}$.

Now $K\cap \overline{V}_{P_i}$ is closed in the compact set $K$, hence compact in $M$.
By continuity of $\varphi_{P_i}$ on $U_{P_i}$, the image
\[
E_i=\varphi_{P_i}\bigl(K\cap \overline{V}_{P_i}\bigr)
\]
is compact in $\R^{2n_{\mathrm v}}$. Moreover, from $\overline{V}_{P_i}\subset \varphi_{P_i}^{-1}(B_i)$ we obtain
$E_i\subset B_i$, so $\bigcup_{i=1}^m E_i$ is bounded.

Finally, $B_i$ is convex (a closed ball of a norm), hence $\operatorname{conv}(E_i)\subset B_i$.
Together with $B_i\subset \varphi_{P_i}(U_{P_i})$, this yields
\[
\operatorname{conv}(E_i)\subset \varphi_{P_i}(U_{P_i}),\qquad i=1,\dots,m.
\]
\end{proof}

\section{Proof for Lemma~\ref{lem:volumetric}}
\begin{proof}
Since $E\subset B_2(0,R)$, monotonicity of covering numbers gives
$\mathcal{N}_\delta(E;\|\cdot\|_2)\le \mathcal{N}_\delta(B_2(0,R);\|\cdot\|_2)$.

If $\delta\ge R$, then $B_2(0,R)\subset B_2(0,\delta)$, hence
$\mathcal{N}_\delta(B_2(0,R);\|\cdot\|_2)=1$.

Assume now $0<\delta<R$. Let $\{x_j\}_{j=1}^n\subset B_2(0,R)$ be a maximal $\delta$-separated set
(i.e., $\|x_i-x_j\|_2>\delta$ for $i\neq j$ and it is maximal by inclusion).
By maximality, $\{x_j\}$ is a $\delta$-net of $B_2(0,R)$, hence
$N_\delta(B_2(0,R);\|\cdot\|_2)\le n$.
Moreover, the balls $\{B_2(x_j,\delta/2)\}_{j=1}^n$ are disjoint and satisfy
$B_2(x_j,\delta/2)\subset B_2(0,R+\delta/2)$, so a volume comparison yields
\[
n\,\mathrm{vol}\big(B_2(0,\delta/2)\big)
\ \le\ \mathrm{vol}\big(B_2(0,R+\delta/2)\big)
\ =\ \Big(\frac{R+\delta/2}{\delta/2}\Big)^{2n_{\mathrm v}}\mathrm{vol}\big(B_2(0,\delta/2)\big).
\]
Therefore $n\le (1+2R/\delta)^{2n_{\mathrm v}}$. Since $\delta<R$ implies $1+2R/\delta\le 3R/\delta$, we obtain
$\mathcal{N}_\delta(B_2(0,R);\|\cdot\|_2)\le (3R/\delta)^{2n_{\mathrm v}}$.
Combining the cases $\delta\ge R$ and $0<\delta<R$ proves the claim for some absolute constant
$C_{\mathrm{vol}}$ (e.g.\ one may take $C_{\mathrm{vol}}=3$).
\end{proof}

\section{Effect of the $1/\sqrt{N}$}
\label{sec:normalization}
Although the TV bound \Eqref{eq:TV_bound_cond} does not display an explicit dependence on the discretization
level $N$ (or the ambient dimension $d_N$), this is facilitated by the $1/\sqrt{N}$ normalization built into
the measurement map $F_N$.

To illustrate its role, consider the unnormalized measurements
\[
\widetilde F_N(\gamma)\ \coloneqq \ \sqrt{N}\,F_N(\gamma),\qquad \widetilde F_N(K)\ \coloneqq \ \sqrt{N}\,F_N(K).
\]
For every $\varepsilon>0$, scaling implies
\[
\mathcal{N}_\varepsilon\!\big(\widetilde F_N(K);\|\cdot\|_2\big)\ =\ \mathcal{N}_{\varepsilon/\sqrt{N}}\!\big(F_N(K);\|\cdot\|_2\big).
\]
Using the covering number bound from the proof of Theorem~\ref{thm:FN-main}, there exists a constant
$C>0$ (independent of $N$) such that
\[
\mathcal{N}_\eta\big(F_N(K)\big)\ \le\ m\max\Big\{1,\Big(\frac{C}{\eta}\Big)^{2n_{\mathrm v}}\Big\}.
\]
Taking $\varepsilon=\epsilon_0=T^{-c_{\epsilon_0}}$ yields an additional additive $n_{\mathrm v}\log N$ term in
$\log \mathcal{N}_{\epsilon_0}(\widetilde F_N(K))$ (corresponding to the factor $(\sqrt{N})^{2n_{\mathrm v}}=N^{n_{\mathrm v}}$),
while the intrinsic exponent $2n_{\mathrm v}$ is unchanged.
Similarly, the support radius rescales as
$\sup_{\gamma\in K}\|\widetilde F_N(\gamma)\|_2=\sqrt{N}\sup_{\gamma\in K}\|F_N(\gamma)\|_2\le \sqrt{N}\,r_K$.

\end{document}
